\documentclass[10pt,a4paper]{article}

\usepackage[utf8]{inputenc}
\usepackage[T1]{fontenc}
\usepackage{amsmath,mathtools,amsfonts,amssymb}
\usepackage{graphicx}
\usepackage[spanish, es-tabla]{babel}
\usepackage{lastpage,tabto,xcolor}
\usepackage[a4paper, top=0.9in, bottom=1in, left=0.9in, right=0.9in]{geometry}
\usepackage{fancyhdr}
\usepackage{titling}
\usepackage{float}
\usepackage{enumitem}
\usepackage{booktabs}
\usepackage[mode=image]{standalone}

\usepackage{tikz}
\usepackage{pgfplots}
\pgfplotsset{compat=1.18}
\usepackage{caption}
\usepackage{framed}
\usetikzlibrary{arrows.meta, babel}

\definecolor{utnblue}{RGB}{0, 51, 102}
\definecolor{utnred}{RGB}{204, 0, 0}
\definecolor{utngreen}{RGB}{0, 153, 76}

% Fecha de versión automática según fecha de compilación
\newcommand{\verdate}{\the\year.\ifnum\month<10 0\fi\the\month.\ifnum\day<10 0\fi\the\day}

% Incisos con letras: a), b), c), ...
\setlist[enumerate,1]{label=\textbf{\alph*)}, leftmargin=*, itemsep=4pt}

\setlength{\headsep}{12pt}
\setlength{\droptitle}{-0.5in}
\pretitle{\begin{center}\vspace{-1em}\normalfont\Large}
\posttitle{\par\end{center}\vspace{-2.5em}}
\preauthor{\begin{center}\vspace{-0.5em}}
\postauthor{\par\end{center}}
\predate{}\postdate{}

\pagestyle{fancy}
\lhead{\footnotesize Análisis de Señales y Sistemas  - Año \the\year{} Ver. \verdate}
\rhead{\footnotesize Resolución del Trabajo Práctico N$^\text{o}$6}
\rfoot{\footnotesize Página \thepage\ de \pageref{LastPage}}
\renewcommand{\headrulewidth}{0.4pt}
\renewcommand{\footrulewidth}{0.4pt}

% Encabezado de tema
\newcommand{\tema}[1]{%
  \bigskip
  \section*{#1}
  \vspace{-0.6em}
}

% Encabezado de ejercicio
\newcommand{\ejercicio}[1]{\bigskip\noindent{\large\textbf{Ejercicio #1}}\par\nobreak\smallskip}

% Consigna transcripta del enunciado
\newenvironment{consigna}{%
  \par\smallskip
  \def\FrameCommand{\textcolor{utnblue!35}{\vrule width 2pt}\hspace{10pt}}%
  \MakeFramed{\advance\hsize-\width \FrameRestore}%
  \small\noindent\textit{Enunciado.}\enspace\ignorespaces
}{\endMakeFramed\par\smallskip}

\title{%
	{\normalsize Departamento de Ingeniería en Tecnologías Electrónicas, UTN FRM}\\[0.1cm]
	{\large Análisis de Señales y Sistemas}\\[0.15cm]
	{\normalsize Resolución del Trabajo Práctico N$^{\text o}$6: Señales de tiempo discreto y teorema del muestreo}%
}
\author{}
\date{}

\begin{document}
	\maketitle
	\thispagestyle{fancy}

	\label{sec:resolucion_tp6}

	\noindent
	Antes de cada bloque de ejercicios va un repaso breve de las herramientas que se usan
	en él. La notación es la que se viene usando: $\Delta t$ es el intervalo de muestreo,
	$\omega_s = 2\pi/\Delta t$ la frecuencia de muestreo, $\lambda = \omega\,\Delta t$ la
	frecuencia digital ---un ángulo, medido en radianes por muestra---, la unidad
	imaginaria es $j$ y los ángulos se expresan en radianes con su equivalencia en grados
	entre paréntesis. Las señales de tiempo continuo llevan paréntesis, $x(t)$; las
	secuencias, corchetes, $x[n]$.

	% ============================================================================
	\tema{Del mundo continuo al digital}
	% ============================================================================

	La conversión analógico-digital son dos operaciones independientes, sobre ejes
	distintos. El \textbf{muestreo} actúa sobre el tiempo y produce la secuencia
	\[
		x[n] = x(n\,\Delta t),
	\]
	donde $n$ cuenta muestras, no segundos. La \textbf{cuantización} actúa sobre la
	amplitud: redondea cada valor a la grilla de niveles que la palabra binaria puede
	nombrar. Con $B$ bits hay $2^B$ palabras distintas.

	En el formato \emph{signo-magnitud} de $B$ bits, uno de los bits carga el signo y los
	$B-1$ restantes la magnitud, que recorre los enteros de $0$ a $2^{B-1}-1$. Si el paso
	de la grilla es $q$, el fondo de escala alcanza $\pm(2^{B-1}-1)\,q$, y el error de
	redondeo nunca supera medio paso:
	\[
		|e| \leq \frac{q}{2}.
	\]

	% ----------------------------------------------------------------------------
	\ejercicio{1: muestreo y cuantización de una senoide}

	\begin{consigna}
	Una tensión senoidal $x(t) = 1{,}2\cos(2\pi\cdot 100\,t)$ V se
	muestrea cada $\Delta t = 1{,}25$ ms y se cuantiza después con cuatro
	bits por muestra en signo-magnitud, con un paso de grilla de $0{,}2$ V.
	\begin{enumerate}
		\item Obtener $x[n]$ y simplificar el argumento del coseno hasta
		dejarlo en la forma $\cos(\lambda n)$.
		\item Indicar cuántas muestras entran en un ciclo de $x(t)$ y cuál
		es el período fundamental $N$ de la secuencia.
		\item Verificar que la grilla de niveles cubre todo el rango de la
		señal, e indicar cuántos niveles distintos quedan definidos.
		\item Completar una tabla con la muestra exacta, la magnitud
		entera, la palabra de cuatro bits y el valor cuantizado, para
		$n = 0, 1, \dots, 7$.
		\item Indicar el error máximo de redondeo.
	\end{enumerate}
	\end{consigna}

	\begin{enumerate}
		\item Muestrear es evaluar la señal en $t = n\,\Delta t$:
		\[
			x[n] = x(n\,\Delta t)
			     = 1{,}2\cos\!\big(2\pi\cdot 100 \cdot n\cdot 1{,}25\times 10^{-3}\big)
			     = 1{,}2\cos(0{,}25\pi\, n).
		\]
		El coeficiente que acompaña a $n$ dentro del coseno es, por definición, la frecuencia
		digital:
		\[
			\boxed{\;x[n] = 1{,}2\cos\!\Big(\frac{\pi}{4}n\Big),
			\qquad \lambda = \omega\,\Delta t = 2\pi\cdot 100\cdot 1{,}25\times10^{-3}
			     = \frac{\pi}{4}\ \text{rad/muestra}\;(45^\circ).\;}
		\]
		Ese número es el ángulo que el vector giratorio asociado avanza entre una muestra y
		la siguiente, $\pi/4$ rad ($45^\circ$).

		\item Si cada muestra avanza $\pi/4$,
		completar la vuelta de $2\pi$ requiere
		\[
			\frac{2\pi}{\lambda} = \frac{2\pi}{\pi/4} = 8 \ \text{muestras por ciclo}.
		\]
		El mismo número sale del lado físico: el período de $x(t)$ es
		$T_0 = 1/100 = 10$ ms, y $T_0/\Delta t = 10/1{,}25 = 8$.

		Para el período fundamental de la secuencia hay que verificar que la cuenta cierre con
		enteros. La condición de periodicidad pide $\lambda/2\pi = k/N$ racional:
		\[
			\frac{\lambda}{2\pi} = \frac{\pi/4}{2\pi} = \frac{1}{8},
		\]
		fracción ya irreducible, así que $N = 8$ y $k = 1$. La secuencia se repite cada ocho
		muestras y en esas ocho muestras el vector da exactamente una vuelta. Aquí las dos
		cuentas coinciden ---ocho muestras por ciclo y $N = 8$--- porque $k=1$; más adelante
		veremos casos donde no coinciden.

		\item Con cuatro bits en signo-magnitud, un bit lleva el signo y
		quedan tres para la magnitud, que recorre los enteros $m = 0, 1, \dots, 7$. Con paso
		$q = 0{,}2$ V, el mayor valor representable es
		\[
			7 \times 0{,}2 = 1{,}4\ \text{V} \;>\; 1{,}2\ \text{V},
		\]
		de modo que la grilla cubre todo el rango de la señal, que va de $-1{,}2$ a
		$+1{,}2$ V, y sobra un escalón en cada extremo.

		Sobre la cantidad de niveles hay una sutileza que conviene no pasar por alto. Las
		cuatro posiciones binarias producen $2^4 = 16$ palabras distintas, pero \texttt{0000}
		y \texttt{1000} nombran ambas al cero ---``$+0$'' y ``$-0$''---, así que los niveles
		de tensión efectivamente distintos son
		\[
			\boxed{\;2^4 - 1 = 15 \ \text{niveles}\;}
		\]
		repartidos entre $-1{,}4$ V y $+1{,}4$ V.

		\item El procedimiento, muestra por muestra, es: evaluar $x[n]$,
		dividir por el paso, redondear al entero más cercano ---esa es la magnitud $m$---,
		escribir el signo en el bit más significativo y reconstruir el valor cuantizado como
		$x_q[n] = \pm m\,q$. Sobre $n=1$, por ejemplo:
		\[
			x[1] = 1{,}2\cos(\pi/4) = 0{,}8485\ \text{V},
			\qquad \frac{0{,}8485}{0{,}2} = 4{,}24 \;\to\; m = 4,
			\qquad \texttt{0100} \;\to\; x_q[1] = 0{,}8\ \text{V}.
		\]
		Repitiendo para las ocho muestras del período queda la Tabla~\ref{tab:sol_ej1}.

		La novena muestra repetiría la primera, como corresponde a $N=8$.

		\item Redondear al nivel más cercano no puede apartarse
		más de medio paso de grilla:
		\[
			|e| \leq \frac{q}{2} = \frac{0{,}2}{2} = 0{,}1\ \text{V}.
		\]
		La tabla lo confirma en el peor caso: $|0{,}8485 - 0{,}8| = 0{,}0485$ V.
	\end{enumerate}

	\begin{table}[H]
		\centering
		\caption{Muestreo y cuantización de $x(t)=1{,}2\cos(2\pi\cdot100\,t)$ con
		$\Delta t = 1{,}25$ ms y cuatro bits en signo-magnitud, paso $0{,}2$ V.}
		\label{tab:sol_ej1}
		\begin{tabular}{ccccc}
			\toprule
			$n$ & $x[n]$ [V] & magnitud $m$ & palabra & $x_q[n]$ [V] \\
			\midrule
			0 & $ 1{,}2000$ & 6 & \texttt{0110} & $ 1{,}2$ \\
			1 & $ 0{,}8485$ & 4 & \texttt{0100} & $ 0{,}8$ \\
			2 & $ 0{,}0000$ & 0 & \texttt{0000} & $ 0{,}0$ \\
			3 & $-0{,}8485$ & 4 & \texttt{1100} & $-0{,}8$ \\
			4 & $-1{,}2000$ & 6 & \texttt{1110} & $-1{,}2$ \\
			5 & $-0{,}8485$ & 4 & \texttt{1100} & $-0{,}8$ \\
			6 & $ 0{,}0000$ & 0 & \texttt{0000} & $ 0{,}0$ \\
			7 & $ 0{,}8485$ & 4 & \texttt{0100} & $ 0{,}8$ \\
			\bottomrule
		\end{tabular}
	\end{table}

	% ============================================================================
	\tema{La secuencia: anatomía y operaciones}
	% ============================================================================

	Las transformaciones de la variable independiente son las del tiempo continuo, con una
	restricción que lo cambia todo: la secuencia solo existe sobre los enteros.

	\begin{itemize}[nosep]
		\item \textbf{Desplazamiento} $x[n-n_0]$, con $n_0$ \emph{entero}: corre la
		      secuencia $n_0$ lugares a la derecha si $n_0>0$.
		\item \textbf{Rebatimiento} $x[-n]$: refleja alrededor de $n=0$.
		\item \textbf{Decimación} $y[n]=x[Mn]$: conserva una de cada $M$ muestras y
		      descarta el resto. Es \emph{irreversible}.
		\item \textbf{Interpolación} $w[n]=x[n/L]$ si $n$ es múltiplo de $L$, y $0$ en
		      otro caso: aleja cada muestra al múltiplo de su posición y deja $L-1$ ceros
		      entre muestras originales. Es \emph{reversible}.
		\item \textbf{Descomposición par-impar}:
		      $x_p[n] = \tfrac12\big(x[n]+x[-n]\big)$,
		      $x_i[n] = \tfrac12\big(x[n]-x[-n]\big)$.
	\end{itemize}

	% ----------------------------------------------------------------------------
	\ejercicio{2: transformaciones de la variable independiente}

	\begin{consigna}
	Sea la secuencia
	\[
		x[n] = \{\, -1 \quad 2 \quad \underset{\uparrow}{3} \quad 1 \quad -2 \,\}.
	\]
	\begin{enumerate}
		\item Graficar $x[n]$ e indicar su soporte.
		\item Graficar $x[n-3]$, $x[-n]$ y $x[2-n]$, indicando en cada caso
		el soporte resultante.
		\item Indicar cuáles de las cuatro secuencias anteriores son
		causales.
	\end{enumerate}
	\end{consigna}

	La notación por extensión con la flecha en $n=0$ dice que
	\[
		x[-2]=-1,\quad x[-1]=2,\quad x[0]=3,\quad x[1]=1,\quad x[2]=-2,
	\]
	y que fuera de esos índices la secuencia vale cero.

	\begin{enumerate}
		\item El soporte ---el intervalo de índices donde
		la secuencia no es nula--- es $[-2,2]$. La gráfica es el panel (a) de la
		Figura~\ref{fig:sol_ej2}: un segmento vertical por muestra, rematado con un punto.

		\item El soporte se razona antes de mover un solo
		valor, porque así se sabe de antemano dónde tiene que caer cada muestra.

		\emph{Desplazamiento $x[n-3]$.} La muestra que estaba en el índice $m$ pasa al índice
		$m+3$: el soporte $[-2,2]$ se desplaza a $[1,5]$. Verificando un valor,
		$x[3-3]=x[0]=3$, así que el $3$ queda en $n=3$.

		\emph{Rebatimiento $x[-n]$.} La muestra del índice $m$ pasa al índice $-m$. El soporte
		$[-2,2]$ es simétrico, así que no se mueve; lo que se invierte es el orden de
		los valores: $\{\,-2\ \ 1\ \ \underset{\uparrow}{3}\ \ 2\ \ -1\,\}$.

		\emph{Rebatimiento y desplazamiento $x[2-n]$.} Aquí hay que tener cuidado con el
		orden, y por eso vale recorrer los dos caminos posibles.

		\emph{Primer camino: rebatir y después desplazar.} Se define $g[n]=x[-n]$ y se le
		aplica un desplazamiento de dos lugares a la derecha:
		\[
			g[n-2] = x\big[-(n-2)\big] = x[2-n]. \quad\checkmark
		\]
		\emph{Segundo camino: desplazar y después rebatir.} Se define $h[n]=x[n+2]$
		---un \emph{adelanto}, no un retardo--- y se lo rebate:
		\[
			h[-n] = x[-n+2] = x[2-n]. \quad\checkmark
		\]
		Los dos llegan al mismo lugar, pero por rutas distintas: en el primero el
		desplazamiento es a la derecha y ocurre \emph{después} del rebatimiento; en el segundo
		es a la izquierda y ocurre \emph{antes}. El error habitual es mezclarlos ---desplazar a
		la derecha primero y rebatir después---, que produce $x[-n-2]$, otra secuencia.

		El soporte se obtiene de la condición $-2 \leq 2-n \leq 2$, es decir
		$0 \leq n \leq 4$, y los valores quedan
		$\{\,\underset{\uparrow}{-2}\ \ 1\ \ 3\ \ 2\ \ -1\,\}$. Las cuatro secuencias están en
		la Figura~\ref{fig:sol_ej2}.

		\item Una secuencia es causal si vale cero para todo
		$n<0$. Mirando los soportes: $x[n]$ no lo es ($x[-2]\neq0$) y $x[-n]$ tampoco. En
		cambio $x[n-3]$, con soporte $[1,5]$, y $x[2-n]$, con soporte
		$[0,4]$, sí lo son. Vale la pena notar el mecanismo: el desplazamiento a la
		derecha es lo que empuja el soporte al semieje positivo, y por eso dos de las cuatro
		resultan causales.
	\end{enumerate}

	\begin{figure}[H]
		\centering
		\includestandalone[width=0.56\linewidth]{figures/fig_sol_ej2/fig_sol_ej2}
		\caption{Ejercicio~2. (a) $x[n]$, soporte $[-2,2]$.
		(b) $x[n-3]$, soporte $[1,5]$. (c) $x[-n]$, soporte $[-2,2]$.
		(d) $x[2-n]$, soporte $[0,4]$.}
		\label{fig:sol_ej2}
	\end{figure}

	% ----------------------------------------------------------------------------
	\ejercicio{3: decimación e interpolación}

	\begin{consigna}
	Sea la secuencia
	\[
		x[n] = \{\, \underset{\uparrow}{1} \quad -1 \quad 2 \quad 0 \quad 1 \quad 3 \,\}.
	\]
	\begin{enumerate}
		\item Determinar y graficar la decimación $y[n] = x[2n]$.
		\item Determinar y graficar la interpolación $w[n] = x[n/3]$.
		\item Indicar si se puede recuperar $x[n]$ a partir de $y[n]$, y si
		se puede a partir de $w[n]$. Justificar en los dos casos.
		\item Explicar por qué el escalado temporal $x(at)$ del tiempo continuo no
		se traslada al discreto como una sola operación.
	\end{enumerate}
	\end{consigna}

	La secuencia de partida es causal, con $x[0]=1$, $x[1]=-1$, $x[2]=2$, $x[3]=0$,
	$x[4]=1$, $x[5]=3$.

	\begin{enumerate}
		\item Se evalúa $x$ en los índices pares:
		\[
			y[0]=x[0]=1,\quad y[1]=x[2]=2,\quad y[2]=x[4]=1,\quad y[3]=x[6]=0,
		\]
		y hacia la izquierda $y[-1]=x[-2]=0$. Entonces
		\[
			\boxed{\;y[n]=\{\,\underset{\uparrow}{1}\quad 2\quad 1\,\}\;}
		\]
		De las seis muestras originales sobreviven tres: las de índice par. Las de índice
		impar ---los valores $-1$, $0$ y $3$--- desaparecen.

		\item La definición pide que $n/3$ sea entero;
		donde no lo es, la secuencia vale cero. Las muestras heredadas de $x$ caen en los
		múltiplos de tres:
		\[
			w[0]=x[0]=1,\ \ w[3]=x[1]=-1,\ \ w[6]=x[2]=2,\ \ w[9]=x[3]=0,\ \
			w[12]=x[4]=1,\ \ w[15]=x[5]=3,
		\]
		y entre cada par de ellas quedan dos ceros ---en general $L-1$, con $L=3$---. La
		Figura~\ref{fig:sol_ej3} muestra las tres secuencias sobre el mismo eje, de modo que
		se vea cómo una se comprime y la otra se estira.

		\item A partir de $y[n]$ \emph{no} se puede recuperar $x[n]$.
		La decimación descartó $x[1]$, $x[3]$ y $x[5]$, y esos valores no están en ninguna
		parte de $y$: cualquier secuencia que coincida con $x$ en los índices pares produce la
		misma $y$. La operación borra información.

		A partir de $w[n]$ sí se puede, y la fórmula es inmediata: $x[n] = w[3n]$. La
		interpolación no perdió nada ---todas las muestras originales siguen ahí, solo que
		separadas por ceros---, así que decimar por el mismo factor la deshace.

		La asimetría tiene una lectura simple: la decimación descarta muestras, la interpolación
		agrega ceros. Solo la segunda es reversible.

		\item En tiempo continuo, $x(at)$ es una
		sola operación bien definida para cualquier $a$ real: la señal existe en todo instante,
		así que $x(at)$ tiene un valor para cada $t$. En tiempo discreto, la secuencia solo
		tiene valores sobre los enteros, y ahí las dos direcciones se comportan de manera distinta:
		\begin{itemize}[nosep]
			\item Si $|a|>1$, por ejemplo $x[2n]$, el índice $2n$ siempre es entero y la operación
			      está definida, pero se saltea muestras: es una decimación, y pierde
			      información.
			\item Si $|a|<1$, por ejemplo $x[n/3]$, el índice $n/3$ deja de ser entero para casi
			      todo $n$, y ahí $x$ sencillamente no está definida. Hay que decidir qué
			      poner en esas posiciones ---la interpolación pone ceros---.
		\end{itemize}
		Por eso lo que en el continuo es una única operación, en el discreto se parte en dos
		operaciones distintas, con propiedades opuestas. Rellenar esos ceros con valores
		razonables ya no es una operación sobre el índice sino tarea de un sistema.
	\end{enumerate}

	\begin{figure}[H]
		\centering
		\includestandalone[width=0.78\linewidth]{figures/fig_sol_ej3/fig_sol_ej3}
		\caption{Ejercicio~3. (a) $x[n]$. (b) Decimación $y[n]=x[2n]$: quedan las
		muestras de índice par. (c) Interpolación $w[n]=x[n/3]$: cada muestra se aleja
		al triple de su posición y entre ellas aparecen dos ceros.}
		\label{fig:sol_ej3}
	\end{figure}

	% ----------------------------------------------------------------------------
	\ejercicio{4: componentes par e impar}

	\begin{consigna}
	Dada la secuencia
	\[
		x[n] = \{\, 6 \quad 4 \quad \underset{\uparrow}{-2} \quad -2 \,\},
	\]
	\begin{enumerate}
		\item Determinar sus componentes par e impar, volcando el cálculo
		en una tabla índice por índice.
		\item Graficar las tres secuencias.
		\item Verificar que $x_i[0] = 0$ y que la suma de las componentes
		reconstruye $x[n]$.
	\end{enumerate}
	\end{consigna}

	La secuencia es $x[-2]=6$, $x[-1]=4$, $x[0]=-2$, $x[1]=-2$, y cero fuera.

	\begin{enumerate}
		\item El primer paso es siempre construir $x[-n]$,
		porque las dos fórmulas la necesitan. Rebatir intercambia el valor del índice $m$
		con el del índice $-m$:
		\[
			x[-n]\big|_{n=-2}=x[2]=0,\quad
			x[-n]\big|_{n=-1}=x[1]=-2,\quad
			x[-n]\big|_{n=1}=x[-1]=4,\quad
			x[-n]\big|_{n=2}=x[-2]=6.
		\]
		Con eso se aplican las dos fórmulas índice por índice. En $n=-1$, por ejemplo,
		\[
			x_p[-1] = \frac{x[-1]+x[1]}{2} = \frac{4+(-2)}{2} = 1,
			\qquad
			x_i[-1] = \frac{x[-1]-x[1]}{2} = \frac{4-(-2)}{2} = 3.
		\]
		La Tabla~\ref{tab:sol_ej4} recoge el cálculo completo.

		\item Las tres secuencias están en la
		Figura~\ref{fig:sol_ej4}. Se ve de un vistazo lo que la tabla dice con números:
		$x_p$ es simétrica respecto del eje vertical y $x_i$ es antisimétrica.

		\item La primera es que $x_i[0]=0$, y así
		resulta, porque la fórmula en $n=0$ pide $\tfrac12(x[0]-x[0])$, que vale cero
		cualquiera sea la secuencia. Toda componente impar se anula en el origen.

		La segunda es que las componentes reconstruyan la señal. Sumando las dos últimas
		filas de la tabla, columna por columna,
		\[
			3+3=6,\qquad 1+3=4,\qquad -2+0=-2,\qquad 1+(-3)=-2,\qquad 3+(-3)=0,
		\]
		que es exactamente la fila de $x[n]$. La razón es algebraica y no depende del
		ejemplo:
		\[
			x_p[n]+x_i[n] = \frac{x[n]+x[-n]}{2}+\frac{x[n]-x[-n]}{2} = x[n].
		\]
	\end{enumerate}

	\begin{table}[H]
		\centering
		\caption{Descomposición par-impar índice por índice.}
		\label{tab:sol_ej4}
		\begin{tabular}{lrrrrr}
			\toprule
			$n$      & $-2$ & $-1$ & $0$  & $1$  & $2$ \\
			\midrule
			$x[n]$   & $6$  & $4$  & $-2$ & $-2$ & $0$ \\
			$x[-n]$  & $0$  & $-2$ & $-2$ & $4$  & $6$ \\
			$x_p[n]$ & $3$  & $1$  & $-2$ & $1$  & $3$ \\
			$x_i[n]$ & $3$  & $3$  & $0$  & $-3$ & $-3$ \\
			\bottomrule
		\end{tabular}
	\end{table}

	\begin{figure}[H]
		\centering
		\includestandalone[width=0.6\linewidth]{figures/fig_sol_ej4/fig_sol_ej4}
		\caption{Ejercicio~4. (a) $x[n]$. (b) Componente par $x_p[n]$, simétrica
		respecto de $n=0$. (c) Componente impar $x_i[n]$, antisimétrica y nula en
		$n=0$.}
		\label{fig:sol_ej4}
	\end{figure}

	Vale notar cómo cambian los soportes. El de $x$ es $[-2,1]$, y los de $x_p$ y $x_i$ son
	$[-2,2]$: un índice más ancho, porque rebatir arrastra las muestras de un lado al otro y
	las dos componentes heredan las dos alas. En $x_i$ hay que agregar una salvedad: aunque el
	origen queda dentro del soporte, ahí la componente impar se anula, porque $x_i[0]=0$
	cualquiera sea la secuencia.

	% ============================================================================
	\tema{La exponencial compleja discreta y la frecuencia digital}
	% ============================================================================

	La exponencial compleja discreta $e^{j\lambda n}$ sale de fotografiar el vector
	giratorio $e^{j\omega t}$ cada $\Delta t$ segundos:
	\[
		e^{j\omega(n\Delta t)} = e^{j(\omega\Delta t)n} = e^{j\lambda n},
		\qquad \lambda = \omega\,\Delta t .
	\]
	La frecuencia digital $\lambda$ \emph{es un ángulo}: el que avanza el vector entre foto
	y foto. De ahí salen los dos hechos que gobiernan todo este bloque.

	\begin{itemize}[nosep]
		\item \textbf{Repetición cada $2\pi$}: $e^{j(\lambda+2\pi)n} = e^{j\lambda n}$.
		      Todas las frecuencias digitales distintas caben en un intervalo de longitud
		      $2\pi$, que tomamos como $(-\pi,\pi]$.
		\item \textbf{Tope de oscilación en $\lambda=\pi$}: media vuelta por muestra es el
		      cambio más brusco posible; la secuencia alterna $+1,-1,+1,\dots$ Pasado
		      $\pi$, la oscilación se \emph{desacelera}.
	\end{itemize}

	Reducir una frecuencia al intervalo $(-\pi,\pi]$ es restarle el múltiplo de $2\pi$ que
	haga falta. Un $\lambda$ negativo describe un vector que gira hacia atrás.

	% ----------------------------------------------------------------------------
	\ejercicio{5: frecuencias digitales equivalentes}

	\begin{consigna}
	Para cada una de las frecuencias digitales siguientes, hallar la
	frecuencia equivalente en el intervalo $(-\pi,\pi]$:
	\begin{center}
	$\lambda_1 = \dfrac{\pi}{3}$, \quad
	$\lambda_2 = \dfrac{5\pi}{3}$, \quad
	$\lambda_3 = \dfrac{7\pi}{3}$, \quad
	$\lambda_4 = \pi$, \quad
	$\lambda_5 = \dfrac{11\pi}{6}$, \quad
	$\lambda_6 = 3\pi$.
	\end{center}
	\begin{enumerate}
		\item Ordenarlas de menor a mayor velocidad de oscilación.
		\item Indicar qué pares producen exactamente la misma secuencia
		$e^{j\lambda n}$ y qué pares comparten solo la parte real.
		\item Justificar, con la cuenta, por qué sumarle $2\pi$ a $\lambda$
		deja la secuencia intacta.
	\end{enumerate}
	\end{consigna}

	Se le resta a cada frecuencia el múltiplo de $2\pi$ que la deja dentro de
	$(-\pi,\pi]$:
	\[
		\begin{aligned}
			\lambda_1 &= \tfrac{\pi}{3} && \text{ya está en el intervalo} &&\to\ \tfrac{\pi}{3}\ (60^\circ),\\
			\lambda_2 &= \tfrac{5\pi}{3} && \tfrac{5\pi}{3}-2\pi = -\tfrac{\pi}{3} &&\to\ -\tfrac{\pi}{3}\ (-60^\circ),\\
			\lambda_3 &= \tfrac{7\pi}{3} && \tfrac{7\pi}{3}-2\pi = \tfrac{\pi}{3} &&\to\ \tfrac{\pi}{3}\ (60^\circ),\\
			\lambda_4 &= \pi && \text{extremo derecho del intervalo} &&\to\ \pi\ (180^\circ),\\
			\lambda_5 &= \tfrac{11\pi}{6} && \tfrac{11\pi}{6}-2\pi = -\tfrac{\pi}{6} &&\to\ -\tfrac{\pi}{6}\ (-30^\circ),\\
			\lambda_6 &= 3\pi && 3\pi-2\pi = \pi &&\to\ \pi\ (180^\circ).
		\end{aligned}
	\]
	Los cuatro puntos distintos que quedan están sobre el círculo unidad en la
	Figura~\ref{fig:sol_ej5}. Es la vista más útil para el resto del ejercicio: cada
	frecuencia es el ángulo que el vector avanza por muestra, y reducir a $(-\pi,\pi]$ es
	elegir, entre todos los ángulos que llevan al mismo punto, el de menor tamaño.

	\begin{figure}[H]
		\centering
		\includestandalone[width=0.48\linewidth]{figures/fig_sol_ej5/fig_sol_ej5}
		\caption{Ejercicio~5. Las seis frecuencias, reducidas al intervalo
		$(-\pi,\pi]$, sobre el círculo unidad: cada vector marca el avance angular por
		muestra. En azul los ángulos positivos (giro antihorario) y en rojo los
		negativos (giro horario).}
		\label{fig:sol_ej5}
	\end{figure}

	\begin{enumerate}
		\item La velocidad la fija el
		\emph{tamaño} del paso angular, es decir $|\lambda|$ una vez reducido, sin importar el
		sentido: un vector que retrocede $30^\circ$ por muestra oscila tan rápido como uno que
		avanza $30^\circ$. Entonces
		\[
			\underbrace{|\lambda_5|}_{\pi/6}
			\;<\;
			\underbrace{|\lambda_1| = |\lambda_2| = |\lambda_3|}_{\pi/3}
			\;<\;
			\underbrace{|\lambda_4| = |\lambda_6|}_{\pi}.
		\]
		Las barras están porque dos de las frecuencias reducidas son negativas,
		$\lambda_2 = -\pi/3$ y $\lambda_5 = -\pi/6$, y sus vectores giran en sentido horario.
		La más lenta es $\lambda_5$ y las más rápidas son $\lambda_4$ y $\lambda_6$, que están
		justo en el tope: alternan $+1,-1,+1,\dots$

		\item Aquí hay dos hechos distintos que se
		confunden con facilidad.

		\emph{Misma secuencia $e^{j\lambda n}$}: pasa cuando las frecuencias reducidas son
		\emph{iguales}. Son los pares $(\lambda_1,\lambda_3)$, ambos $\pi/3$, y
		$(\lambda_4,\lambda_6)$, ambos $\pi$. Coinciden la parte real y la imaginaria: son la
		misma secuencia compleja, muestra por muestra.

		\emph{Solo la parte real}: pasa cuando las frecuencias reducidas son \emph{opuestas}.
		Es el caso de $\lambda_2 = -\pi/3$ frente a $\lambda_1 = \lambda_3 = \pi/3$. En efecto,
		\[
			e^{-j\pi n/3} = \overline{e^{\,j\pi n/3}},
		\]
		y dos conjugados tienen la misma parte real y partes imaginarias opuestas:
		\[
			\cos(-\tfrac{\pi}{3}n) = \cos(\tfrac{\pi}{3}n),
			\qquad
			\sin(-\tfrac{\pi}{3}n) = -\sin(\tfrac{\pi}{3}n).
		\]
		Los dos vectores giran a la misma velocidad en sentidos opuestos: comparten el coseno
		---que no distingue el sentido del giro--- pero son secuencias distintas.

		\item La cuenta es corta:
		\[
			e^{j(\lambda+2\pi)n}
			= e^{j\lambda n}\,e^{j2\pi n}
			= e^{j\lambda n}\,\big(e^{j2\pi}\big)^{n}
			= e^{j\lambda n}\cdot 1^{\,n}
			= e^{j\lambda n}.
		\]
		El factor $e^{j2\pi n}$ vale $1$ porque $n$ es entero: el vector completa un número
		entero de vueltas entre muestra y muestra. Si $n$ pudiera valer, por ejemplo, $1/2$,
		entonces
		$e^{j2\pi\cdot 1/2}=e^{j\pi}=-1\neq 1$ y la propiedad se caería. Esa es exactamente la
		razón por la cual el tiempo continuo no tiene nada parecido: allí cada $\omega$ produce
		una señal distinta.
	\end{enumerate}

	% ----------------------------------------------------------------------------
	\ejercicio{6: fotos de un vector giratorio}

	\begin{consigna}
	Un vector giratorio de frecuencia $\omega = 2\pi\cdot 50$ rad/s se
	muestrea cada $\Delta t$ segundos.
	\begin{enumerate}
		\item Calcular $\lambda$ para $\Delta t = 1$ ms, $6$ ms, $15$ ms y
		$20$ ms.
		\item Para cada caso, determinar la frecuencia digital equivalente
		en $(-\pi,\pi]$ y describir el giro aparente: qué ángulo avanza el
		vector entre foto y foto y, cuando ese giro existe, en qué sentido
		ocurre.
		\item Determinar para qué valores de $\Delta t$ el giro puede
		seguirse sin ambigüedad, y analizar qué pasa en el caso límite.
	\end{enumerate}
	\end{consigna}

	El vector gira a $\omega = 2\pi\cdot 50 = 100\pi$ rad/s, es decir, $50$ vueltas por
	segundo.

	\begin{enumerate}
		\item Se aplica $\lambda = \omega\,\Delta t$:
		\[
			\begin{aligned}
				\Delta t = 1\ \text{ms}: &\quad \lambda = 100\pi\cdot 10^{-3} = 0{,}1\pi = \tfrac{\pi}{10} && (18^\circ),\\
				\Delta t = 6\ \text{ms}: &\quad \lambda = 100\pi\cdot 6\times10^{-3} = 0{,}6\pi = \tfrac{3\pi}{5} && (108^\circ),\\
				\Delta t = 15\ \text{ms}: &\quad \lambda = 100\pi\cdot 15\times10^{-3} = 1{,}5\pi = \tfrac{3\pi}{2} && (270^\circ),\\
				\Delta t = 20\ \text{ms}: &\quad \lambda = 100\pi\cdot 20\times10^{-3} = 2\pi && (360^\circ).
			\end{aligned}
		\]

		\item La reducción a $(-\pi,\pi]$ es la que dice qué ve quien
		solo tiene las fotos, porque de todos los ángulos que llevan al mismo punto la
		secuencia registra el más corto.
		\begin{itemize}[nosep]
			\item $\lambda = \pi/10$ ya está dentro del intervalo. El vector avanza
			      $18^\circ$ por foto en sentido antihorario, y eso es también lo que se ve:
			      giro aparente $=$ giro real. Hacen falta $20$ fotos por vuelta.
			\item $\lambda = 3\pi/5$ también está dentro. Avanza $108^\circ$ por foto,
			      antihorario, y así se ve. Ya pasó el cuarto de vuelta por muestra, pero
			      todavía no llega a media vuelta, que es el límite.
			\item $\lambda = 3\pi/2$ se reduce a $3\pi/2 - 2\pi = -\pi/2$ ($-90^\circ$). El
			      vector avanza en realidad tres cuartos de vuelta hacia adelante, pero entre
			      foto y foto aparece un cuarto de vuelta \emph{antes} de donde estaba: se lo
			      ve girar \textbf{en sentido horario}, $90^\circ$ por foto. Es exactamente la
			      rueda que parece girar al revés.
			\item $\lambda = 2\pi$ se reduce a $0$. El vector da una vuelta entera entre foto y
			      foto y vuelve al mismo punto: la secuencia es constante y el vector aparece
			      \textbf{quieto}. Aquí no hay sentido de giro que indicar, porque el ángulo
			      aparente por foto es cero.
		\end{itemize}
		Los dos últimos casos son el mismo fenómeno con distinta intensidad: la secuencia
		registra dónde quedó el vector, no cuántas vueltas dio para llegar ahí.

		\item El giro es inequívoco mientras el avance
		por foto sea menor que media vuelta, porque en cuanto pasa de $\pi$ el camino corto
		entre dos fotos consecutivas es el de sentido contrario. La condición es
		\[
			|\lambda| < \pi
			\quad\Longleftrightarrow\quad
			100\pi\,\Delta t < \pi
			\quad\Longleftrightarrow\quad
			\boxed{\;\Delta t < 10\ \text{ms}.\;}
		\]
		La condición es una desigualdad estricta, así que \emph{no existe} un mayor
		$\Delta t$ admisible; $10$ ms es el supremo y no se alcanza. En el caso
		límite $\Delta t = 10$ ms exactos queda $\lambda = \pi$: el vector da media vuelta por
		foto y la secuencia no permite decidir si giró $+\pi$ o $-\pi$, porque las dos
		rotaciones dejan el vector en el mismo punto. Ese caso límite tiene nombre propio: la
		condición $\Delta t < 10$ ms equivale a $f_s > 100$ Hz, el doble de los $50$ Hz de la
		señal, que es la condición de Nyquist de la segunda parte del trabajo.
	\end{enumerate}

	% ============================================================================
	\tema{Periodicidad y señales armónicas}
	% ============================================================================

	Una secuencia es periódica si $x[n+N]=x[n]$ con $N$ \textbf{entero}, y el período
	fundamental es el menor $N$ que lo cumple. Para la exponencial compleja, imponer esa
	condición conduce a
	\[
		e^{j\lambda N}=1
		\quad\Longleftrightarrow\quad
		\lambda N = 2\pi k
		\quad\Longleftrightarrow\quad
		\frac{\lambda}{2\pi}=\frac{k}{N},
	\]
	con $k$ y $N$ enteros. Entonces:

	\begin{itemize}[nosep]
		\item $e^{j\lambda n}$ es periódica \textbf{si y solo si} $\lambda/2\pi$ es
		      racional.
		\item Escrita $\lambda/2\pi = k/N$ como fracción \emph{irreducible}, $N$ es el
		      período fundamental y $k$ cuenta cuántas vueltas da el vector en ese período.
		\item La fórmula $T_0 = 2\pi/\omega$ del tiempo continuo es el caso particular
		      $k=1$; en general hay que aumentar $k$ hasta que la cuenta cierre con un entero.
		\item Las \textbf{armónicas} de período $N$ son $\phi_k[n]=e^{j(2\pi/N)kn}$, y solo
		      hay $N$ distintas, porque $\phi_{k+N}[n]=\phi_k[n]$.
	\end{itemize}

	Para una suma de dos secuencias periódicas de períodos $N_1$ y $N_2$, el período de la
	suma es el mínimo común múltiplo $\operatorname{mcm}(N_1,N_2)$.

	% ----------------------------------------------------------------------------
	\ejercicio{7: periodicidad de secuencias}

	\begin{consigna}
	Determinar si las siguientes secuencias son periódicas. Cuando lo
	sean, hallar el período fundamental $N$ e indicar, para las que son una
	sola exponencial o sinusoide, cuántas vueltas $k$ da el vector giratorio
	en ese período.
	\begin{enumerate}
		\item $x[n] = \cos(\pi n/5)$
		\item $x[n] = e^{j3\pi n/4}$
		\item $x[n] = \cos(15 n)$
		\item $x[n] = \cos(5\pi n)$
		\item $x[n] = \cos(\pi n/3) + \sin(\pi n/4)$
	\end{enumerate}
	\end{consigna}

	El procedimiento es siempre el mismo: leer $\lambda$ del argumento, formar
	$\lambda/2\pi$ y ver si es racional.

	\begin{enumerate}
		\item $x[n]=\cos(\pi n/5)$. Aquí $\lambda = \pi/5$ y
		\[
			\frac{\lambda}{2\pi} = \frac{\pi/5}{2\pi} = \frac{1}{10},
		\]
		racional e irreducible. Entonces $N=10$ y $k=1$: la secuencia se repite cada diez
		muestras, y en esas diez el vector da una vuelta.

		\item $x[n]=e^{j3\pi n/4}$. Con $\lambda=3\pi/4$,
		\[
			\frac{\lambda}{2\pi} = \frac{3\pi/4}{2\pi} = \frac{3}{8},
		\]
		irreducible. Entonces $N=8$ y $k=3$. Vale la pena leerlo: hacen falta ocho muestras
		para que la secuencia se repita, y durante esas ocho muestras el vector da tres
		vueltas. Es el caso donde la intuición del tiempo continuo ---una vuelta por
		período--- ya no sirve.

		\item $x[n]=\cos(15n)$. Ahora $\lambda = 15$ y
		\[
			\frac{\lambda}{2\pi} = \frac{15}{2\pi},
		\]
		que es irracional, porque $\pi$ lo es y $15$ es entero. No hay enteros $k$ y $N$ que
		satisfagan la condición: la secuencia \textbf{no es periódica}. El vector nunca vuelve
		a pisar exactamente el punto de partida, aunque las muestras se acerquen tanto como se
		quiera. El contraste con el tiempo continuo es directo, porque allí $\cos(15t)$ sí es
		periódica, con $T_0=2\pi/15$; lo que rompe la periodicidad es la exigencia de que el período sea
		entero.

		\item $x[n]=\cos(5\pi n)$. Con $\lambda=5\pi$,
		\[
			\frac{\lambda}{2\pi} = \frac{5\pi}{2\pi} = \frac{5}{2},
		\]
		fracción irreducible, así que $N=2$ y $k=5$. Es periódica de período dos, y en esas
		dos muestras el vector da cinco vueltas.

		Este caso deja ver algo que conviene tener presente. Reduciendo primero la frecuencia
		al intervalo $(-\pi,\pi]$ queda $5\pi - 4\pi = \pi$, y con $\lambda=\pi$ la cuenta
		arroja $\lambda/2\pi = 1/2$, es decir, el mismo $N=2$ pero con $k=1$. Las dos lecturas
		describen la misma secuencia ---que es $\cos(\pi n) = (-1)^n$---: el período $N$ no
		depende del representante que se elija, pero el número de vueltas $k$ sí, porque las
		vueltas de más no quedan registradas en las muestras.

		\item $x[n]=\cos(\pi n/3)+\sin(\pi n/4)$. Al ser una suma, se estudia cada
		término por separado y después se combinan los períodos.
		\[
			\text{Primero: } \lambda_1=\frac{\pi}{3},\quad \frac{\lambda_1}{2\pi}=\frac{1}{6}
			\ \Rightarrow\ N_1=6;
			\qquad
			\text{segundo: } \lambda_2=\frac{\pi}{4},\quad \frac{\lambda_2}{2\pi}=\frac{1}{8}
			\ \Rightarrow\ N_2=8.
		\]
		Los dos son periódicos, así que la suma también lo es, y se repite cuando los dos se
		repiten a la vez: hace falta el menor entero que sea múltiplo de $6$ y de $8$,
		\[
			N = \operatorname{mcm}(6,8) = 24.
		\]
		En esas $24$ muestras el primer término completa $24/6=4$ ciclos y el segundo
		$24/8=3$. Aquí no tiene sentido hablar de un único $k$, porque no hay un solo vector
		girando sino dos, a velocidades distintas.
	\end{enumerate}

	% ----------------------------------------------------------------------------
	\ejercicio{8: elegir $\Delta t$ para fijar el período}

	\begin{consigna}
	La señal $x(t) = 3\cos(20t)$ se muestrea con un intervalo
	$\Delta t$ elegido de modo que la secuencia $x[n]$ resulte periódica con
	período fundamental $N = 10$.
	\begin{enumerate}
		\item Determinar todos los valores posibles de $\lambda$ y de
		$\Delta t$ dentro de una vuelta.
		\item Para cada solución, indicar cuántos ciclos de $x(t)$
		transcurren mientras la secuencia recorre un período.
		\item Señalar cuáles de esas soluciones muestrean por encima del
		doble de la frecuencia de la señal, y relacionar esa condición con
		la posición de $\lambda$ respecto de $\pi$.
	\end{enumerate}
	\end{consigna}

	La señal es $x(t)=3\cos(20t)$, de modo que $\omega = 20$ rad/s, y al muestrear queda
	\[
		x[n]=3\cos(20\,\Delta t\, n) = 3\cos(\lambda n),
		\qquad \lambda = 20\,\Delta t .
	\]
	Este ejercicio recorre la condición de periodicidad al revés: en vez de calcular $N$ a
	partir de $\lambda$, se fija $N$ y se busca qué $\lambda$ lo produce.

	\begin{enumerate}
		\item Pedir período fundamental $N=10$ es pedir
		\[
			\frac{\lambda}{2\pi}=\frac{k}{10}
			\qquad\text{con } \frac{k}{10} \text{ \emph{irreducible}},
		\]
		y ahí está la clave del ejercicio. Si la fracción se pudiera simplificar, el período
		fundamental sería menor que $10$: por ejemplo $k=2$ produce $2/10=1/5$, que es período
		$5$, no $10$. La irreducibilidad exige que $k$ y $10$ no compartan divisores, es decir
		que $k$ no sea múltiplo de $2$ ni de $5$. Dentro de una vuelta ---es decir,
		$0<\lambda<2\pi$, que corresponde a $1\leq k\leq 9$--- los que sobreviven son
		\[
			k \in \{1,\,3,\,7,\,9\}.
		\]
		Para cada uno, $\lambda = 2\pi k/10 = \pi k/5$ y $\Delta t = \lambda/20 = \pi k/100$
		segundos. La Tabla~\ref{tab:sol_ej8} reúne las cuatro soluciones.

		\item El período de la señal
		continua es $T_0 = 2\pi/20 = \pi/10$ s. Un período de la secuencia dura $N$ muestras,
		es decir
		\[
			N\,\Delta t = 10\cdot\frac{\pi k}{100} = \frac{\pi k}{10}\ \text{segundos},
		\]
		y en ese tiempo la señal continua completa
		\[
			\frac{N\,\Delta t}{T_0} = \frac{\pi k/10}{\pi/10} = k \ \text{ciclos}.
		\]
		El resultado es limpio: \emph{el entero $k$ es la cantidad de ciclos de la señal
		continua que transcurren mientras la secuencia recorre un período}. Con $k=1$ se
		muestrea la señal dentro de un solo ciclo, con diez muestras bien repartidas; con
		$k=9$ esas mismas diez muestras se reparten a lo largo de nueve ciclos ---casi un ciclo
		entero entre muestra y muestra---, de modo que provienen de ciclos distintos aunque
		describan la misma secuencia periódica.

		\item La condición de muestrear por encima del doble de
		la frecuencia de la señal es
		\[
			\omega_s > 2\omega
			\quad\Longleftrightarrow\quad
			\frac{2\pi}{\Delta t} > 40
			\quad\Longleftrightarrow\quad
			\Delta t < \frac{\pi}{20} = 157{,}1\ \text{ms}.
		\]
		Mirando la tabla, la cumplen $k=1$ y $k=3$; no la cumplen $k=7$ ni $k=9$.

		Ahora, la condición se lee mucho mejor en frecuencia digital. Multiplicando los dos
		miembros de $\omega_s>2\omega$ por $\Delta t/2$:
		\[
			\frac{\omega_s\,\Delta t}{2} > \omega\,\Delta t
			\quad\Longleftrightarrow\quad
			\frac{2\pi}{2} > \lambda
			\quad\Longleftrightarrow\quad
			\boxed{\;\lambda < \pi\;}
		\]
		donde se usó $\omega_s\,\Delta t = 2\pi$. Es decir: \emph{la condición de Nyquist es
		exactamente que la frecuencia digital caiga por debajo de $\pi$}. La última columna de
		la tabla lo confirma sin necesidad de volver a los milisegundos. Y encaja con el
		Ejercicio~6: $\lambda<\pi$ es la condición para que el vector avance menos de media
		vuelta por foto, es decir, para que el giro se pueda seguir sin ambigüedad.
	\end{enumerate}

	\begin{table}[H]
		\centering
		\caption{Las cuatro soluciones dentro de una vuelta, con los ciclos de $x(t)$ que
		transcurren en un período de la secuencia y el cumplimiento de Nyquist.}
		\label{tab:sol_ej8}
		\begin{tabular}{ccccc}
			\toprule
			$k$ & $\lambda$ & $\Delta t$ [ms] & ciclos de $x(t)$ & $\lambda<\pi$ \\
			\midrule
			$1$ & $\pi/5\ (36^\circ)$   & $31{,}42$  & $1$ & sí \\
			$3$ & $3\pi/5\ (108^\circ)$ & $94{,}25$  & $3$ & sí \\
			$7$ & $7\pi/5\ (252^\circ)$ & $219{,}9$  & $7$ & no \\
			$9$ & $9\pi/5\ (324^\circ)$ & $282{,}7$  & $9$ & no \\
			\bottomrule
		\end{tabular}
	\end{table}

	% ----------------------------------------------------------------------------
	\ejercicio{9: las armónicas de período $N=6$}

	\begin{consigna}
	Se toma $N = 6$ y el conjunto de armónicas
	$\phi_k[n] = e^{j(2\pi/6)kn}$.
	\begin{enumerate}
		\item Escribir las armónicas distintas y verificar con la cuenta que
		$\phi_6[n] = \phi_0[n]$.
		\item Las armónicas $\phi_4$ y $\phi_5$ tienen frecuencia entre
		$\pi$ y $2\pi$. Hallar para cada una la frecuencia equivalente en
		$(-\pi,0)$, verificar sobre un período que la secuencia es la misma e
		indicar qué relación guarda cada una con otra armónica del conjunto.
		Señalar cuál de las seis oscila más rápido.
		\item Comprobar que $\lambda = 2\pi/5$ y $\lambda = 4\pi/5$ tienen
		el mismo período $N = 5$ pero distinta velocidad de oscilación.
	\end{enumerate}
	\end{consigna}

	\begin{enumerate}
		\item Con $N=6$ la fundamental es
		$\lambda_1 = 2\pi/6 = \pi/3$, y las armónicas son sus múltiplos:
		\[
			\phi_k[n] = e^{j\frac{2\pi}{6}kn} = e^{j\frac{\pi}{3}kn},
			\qquad k = 0,1,2,3,4,5,
		\]
		con frecuencias $\lambda_k = \pi k/3$, es decir
		$0,\ \pi/3,\ 2\pi/3,\ \pi,\ 4\pi/3,\ 5\pi/3$. La armónica $k$ da $k$ vueltas mientras
		la secuencia recorre sus seis muestras.

		Que la lista se cierre en seis se comprueba con la cuenta:
		\[
			\phi_6[n] = e^{j\frac{2\pi}{6}\cdot 6 \cdot n} = e^{j2\pi n}
			          = \big(e^{j2\pi}\big)^{n} = 1^{\,n} = 1 = \phi_0[n],
		\]
		donde otra vez el paso decisivo es que $n$ sea entero. El mismo argumento vale para
		cualquier $k$: $\phi_{k+6}[n]=\phi_k[n]\,e^{j2\pi n}=\phi_k[n]$. Es una diferencia
		fuerte con el tiempo continuo, donde las armónicas $e^{jk\omega_0 t}$ son infinitas y
		todas distintas.

		\item Las dos últimas del
		conjunto pasan de $\pi$,
		\[
			\lambda_4 = \frac{4\pi}{3}\ (240^\circ),
			\qquad
			\lambda_5 = \frac{5\pi}{3}\ (300^\circ),
		\]
		y restarle una vuelta a cada una las lleva al tramo negativo:
		\[
			\frac{4\pi}{3} - 2\pi = -\frac{2\pi}{3}\ (-120^\circ),
			\qquad
			\frac{5\pi}{3} - 2\pi = -\frac{\pi}{3}\ (-60^\circ).
		\]
		Que la secuencia no cambie es la cuenta del inciso anterior aplicada a una sola
		armónica:
		\[
			\phi_4[n] = e^{j\frac{4\pi}{3}n}
			          = e^{j\left(-\frac{2\pi}{3}+2\pi\right)n}
			          = e^{-j\frac{2\pi}{3}n}\,e^{j2\pi n}
			          = e^{-j\frac{2\pi}{3}n}.
		\]
		Sobre un período se comprueba muestra por muestra. La Tabla~\ref{tab:sol_ej9} lista el
		ángulo del vector para las dos frecuencias, reducido a una vuelta: las dos filas
		coinciden.

		\begin{table}[H]
			\centering
			\caption{Ejercicio~9, inciso (b). Ángulo del vector, en grados, para las dos
			escrituras de la misma armónica.}
			\label{tab:sol_ej9}
			\begin{tabular}{lcccccc}
				\toprule
				$n$ & $0$ & $1$ & $2$ & $3$ & $4$ & $5$ \\
				\midrule
				$\lambda = 4\pi/3$   & $0$ & $240$ & $480 \equiv 120$ & $720 \equiv 0$
				                     & $960 \equiv 240$ & $1200 \equiv 120$ \\
				$\lambda = -2\pi/3$  & $0$ & $-120 \equiv 240$ & $-240 \equiv 120$
				                     & $-360 \equiv 0$ & $-480 \equiv 240$ & $-600 \equiv 120$ \\
				\bottomrule
			\end{tabular}
		\end{table}

		Escrita con frecuencia negativa, la armónica queda emparentada con otra del conjunto.
		Como $e^{-j\frac{2\pi}{3}n}$ es el conjugado de $e^{j\frac{2\pi}{3}n}$,
		\[
			\phi_4[n] = \overline{\phi_2[n]},
			\qquad\text{y del mismo modo}\qquad
			\phi_5[n] = \overline{\phi_1[n]} .
		\]
		Las seis armónicas son entonces la constante $\phi_0$, la que alterna $+1,-1$
		---$\phi_3$, con $\lambda = \pi$--- y dos pares conjugados, $(\phi_1,\phi_5)$ y
		$(\phi_2,\phi_4)$. Los dos paneles de la Figura~\ref{fig:sol_ej9} muestran el mismo
		conjunto con las dos escrituras: en el panel (a) los seis ángulos se cuentan siempre
		hacia el mismo lado, de $0$ a $2\pi$; en el (b), los dos que pasaban de $\pi$ se
		cuentan hacia atrás, y quedan como los espejos de $\phi_1$ y $\phi_2$ respecto del eje
		real.

		Ese cambio de escritura corrige una lectura equivocada de la velocidad. Reducidas a
		$(-\pi,\pi]$, las frecuencias del conjunto son $0$, $\pm\pi/3$, $\pm2\pi/3$ y $\pi$, y
		la velocidad de oscilación la fija el tamaño del paso angular. La más rápida es
		entonces $\phi_3$, la del medio de la lista. Aumentar $k$ de $3$ a $5$ no acelera nada:
		el vector empieza a girar en el otro sentido y cada vez más despacio, hasta que en
		$k=6$ vuelve a estar quieto.

		\item Se aplica la condición de
		periodicidad a los dos valores propuestos:
		\[
			\lambda = \frac{2\pi}{5}:\ \ \frac{\lambda}{2\pi}=\frac{1}{5}
			\ \Rightarrow\ N=5,\ k=1;
			\qquad
			\lambda = \frac{4\pi}{5}:\ \ \frac{\lambda}{2\pi}=\frac{2}{5}
			\ \Rightarrow\ N=5,\ k=2.
		\]
		Las dos fracciones son irreducibles y tienen el mismo denominador, así que las dos
		secuencias se repiten cada cinco muestras. Pero la segunda avanza $4\pi/5$
		($144^\circ$) por muestra contra $2\pi/5$ ($72^\circ$) de la primera: oscila al doble
		de velocidad y da dos vueltas donde la otra da una.
	\end{enumerate}

	\begin{figure}[H]
		\centering
		\includestandalone[width=0.82\linewidth]{figures/fig_sol_ej9/fig_sol_ej9}
		\caption{Ejercicio~9. Las seis frecuencias armónicas $\lambda_k = 2\pi k/6$
		sobre el círculo unidad, seis puntos igualmente espaciados. (a) Contadas de $0$
		a $2\pi$, como salen de la definición. (b) Reducidas a $(-\pi,\pi]$: las dos que
		pasaban de $\pi$ se cuentan hacia atrás y quedan como los espejos de $\phi_1$ y
		$\phi_2$ respecto del eje real. La armónica $k=6$ vuelve al punto de $k=0$.}
		\label{fig:sol_ej9}
	\end{figure}

	% ============================================================================
	\tema{La exponencial general: la secuencia $z^n$}
	% ============================================================================

	Escribiendo $z$ en forma polar, $z=re^{j\lambda}$, la potencia se factoriza en las dos
	capacidades ya vistas por separado:
	\[
		z^n = \big(re^{j\lambda}\big)^n = r^n\,e^{j\lambda n}.
	\]
	El módulo $r$ escala y el argumento $\lambda$ rota. Tomando partes real e imaginaria
	---y sacando $r^n$ afuera, que es real---,
	\[
		\operatorname{Re}\{z^n\} = r^n\cos(\lambda n),
		\qquad
		\operatorname{Im}\{z^n\} = r^n\sin(\lambda n),
	\]
	una sinusoide con \textbf{envolvente geométrica} $\pm r^n$: amortiguada si $r<1$,
	creciente si $r>1$, de amplitud constante si $r=1$.

	Para clasificar en energía o potencia se usan
	\[
		E_x = \sum_{n=-\infty}^{\infty} |x[n]|^2,
		\qquad
		P_x = \lim_{N\to\infty}\frac{1}{2N+1}\sum_{n=-N}^{N}|x[n]|^2 .
	\]

	% ----------------------------------------------------------------------------
	\ejercicio{10: la secuencia $z^n u[n]$}

	\begin{consigna}
	Sea $z = r\,e^{j\lambda}$ y la secuencia causal
	$x[n] = z^n\,u[n]$.
	\begin{enumerate}
		\item Para $z = 0{,}9\,e^{j\pi/6}$, escribir
		$\operatorname{Re}\{x[n]\}$, identificar la envolvente y el período
		de la oscilación, y graficar las primeras seis muestras.
		\item Repetir para $z = 1{,}1\,e^{j\pi/6}$ y comparar las dos
		gráficas.
		\item Para $z = -0{,}8$, mostrar que $z^n = (0{,}8)^n(-1)^n$ e
		indicar qué valor de $\lambda$ corresponde a ese punto del plano.
		\item Clasificar las secuencias de los incisos (a) y (b) como de
		energía, de potencia o de ninguna de las dos.
		\item Indicar qué puntos del plano complejo hacen que
		$\operatorname{Re}\{x[n]\}$ sea una sinusoide discreta de amplitud
		constante.
	\end{enumerate}
	\end{consigna}

	\begin{enumerate}
		\item $z=0{,}9\,e^{j\pi/6}$. Aplicando la factorización y quedándose con la
		parte real,
		\[
			\boxed{\;\operatorname{Re}\{x[n]\} = (0{,}9)^n\cos\!\Big(\frac{\pi}{6}n\Big)\,u[n].\;}
		\]
		La \textbf{envolvente} es $\pm(0{,}9)^n$: como $r<1$, decae hacia cero, y la sinusoide
		oscila apretada entre esas dos curvas. Para el \textbf{período de la oscilación} se
		aplica la condición de periodicidad al coseno:
		\[
			\lambda = \frac{\pi}{6}
			\quad\Longrightarrow\quad
			\frac{\lambda}{2\pi} = \frac{1}{12}
			\quad\Longrightarrow\quad
			N = 12,\quad k=1,
		\]
		es decir doce muestras por ciclo. El vocabulario acá pide precisión, porque la
		\emph{secuencia} no es periódica ---la envolvente decae y ningún tramo se repite
		idéntico---; lo que tiene período $12$ es la oscilación que va por dentro. Las primeras
		seis muestras ---media oscilación--- están en el panel (a) de la
		Figura~\ref{fig:sol_ej10}.

		\item $z=1{,}1\,e^{j\pi/6}$. Cambia solo el módulo:
		\[
			\operatorname{Re}\{x[n]\} = (1{,}1)^n\cos\!\Big(\frac{\pi}{6}n\Big)\,u[n].
		\]
		El período de la oscilación es el mismo, $N=12$, porque $\lambda$ no cambió. Lo que
		cambia es la envolvente: ahora $r>1$ y $\pm(1{,}1)^n$ se abre en vez de cerrarse. Las
		dos gráficas tienen exactamente los mismos cruces por cero ---sobre las seis muestras
		graficadas, el de $n=3$---; lo único distinto es que en una las amplitudes se achican y
		en la otra crecen sin cota. Los dos parámetros de $z$ actúan de
		manera independiente: $\lambda$ fija la velocidad de oscilación y $r$ la tasa de
		crecimiento.

		\item $z=-0{,}8$. Un número real negativo, escrito en forma polar, tiene módulo
		$0{,}8$ y argumento $\pi$:
		\[
			z = -0{,}8 = 0{,}8\,e^{j\pi}
			\quad\Longrightarrow\quad
			z^n = (0{,}8)^n\,e^{j\pi n} = (0{,}8)^n\,\big(e^{j\pi}\big)^n = (0{,}8)^n(-1)^n,
		\]
		que es lo que se pedía mostrar. La frecuencia digital es
		\[
			\boxed{\;\lambda = \pi\ \text{rad/muestra}\;(180^\circ),\;}
		\]
		el tope de oscilación: media vuelta por muestra. La secuencia decae en módulo y cambia
		de signo en cada paso, que es el comportamiento más brusco que una secuencia
		amortiguada puede tener. Todo el eje real negativo del plano complejo produce este
		signo alternante.

		\item Para la secuencia del inciso (a), acotando el coseno
		por uno,
		\[
			E_x = \sum_{n=0}^{\infty}\big|(0{,}9)^n\cos(\pi n/6)\big|^2
			    \;\leq\; \sum_{n=0}^{\infty}(0{,}81)^n
			    = \frac{1}{1-0{,}81} = 5{,}26 \;<\; \infty,
		\]
		donde se usó la suma de la serie geométrica de razón $0{,}81<1$. La energía es finita,
		así que es una \textbf{señal de energía} ---y, como toda señal de energía finita, su
		potencia media vale cero---.

		Para la del inciso (b) el argumento es el opuesto. El factor $(1{,}1)^n$ crece sin
		cota, así que los términos de la suma no tienden a cero y la energía diverge. Tampoco
		hay potencia finita: al promediar $|x[n]|^2$ sobre una ventana de $2N+1$ muestras, el
		numerador crece como $(1{,}21)^N$ mientras el denominador crece solo como $N$, de modo
		que el cociente también diverge. Es una secuencia que no es \textbf{ni de energía ni de
		potencia}.

		Vale aclarar por qué el enunciado pide la versión causal $z^n u[n]$: si la secuencia
		estuviera definida para todo entero, el caso (a) tendría problemas del otro lado, ya
		que $(0{,}9)^n$ crece sin cota cuando $n\to-\infty$, y ninguna de las dos sería de
		energía.

		\item La parte real es
		$\operatorname{Re}\{x[n]\} = r^n\cos(\lambda n)\,u[n]$, así que la amplitud la fija la
		envolvente $r^n$, y esa envolvente es constante solo si $r^n=1$ para todo $n$, es decir
		si
		\[
			r = |z| = 1 .
		\]
		Los puntos buscados son entonces los del \textbf{círculo unidad}, $z=e^{j\lambda}$.
		Ahí $\operatorname{Re}\{z^n\}=\cos(\lambda n)$ es una sinusoide discreta pura. Adentro
		del círculo ($r<1$) la sinusoide se amortigua y afuera ($r>1$) crece: la posición del
		punto en el plano determina por completo la forma de la señal.
	\end{enumerate}

	\begin{figure}[H]
		\centering
		\includestandalone[width=0.86\linewidth]{figures/fig_sol_ej10/fig_sol_ej10}
		\caption{Ejercicio~10. Parte real de $z^n u[n]$ con $\lambda=\pi/6$ y
		envolvente $\pm r^n$ punteada, sobre las primeras seis muestras ---media
		oscilación, que tiene período $N=12$---. (a) $r=0{,}9$: la oscilación se
		amortigua. (b) $r=1{,}1$: la oscilación crece. Las dos comparten el cruce por
		cero de $n=3$; solo cambia la envolvente. Las escalas verticales de los dos
		paneles son distintas.}
		\label{fig:sol_ej10}
	\end{figure}

	% ============================================================================
	\tema{Escalón, impulso y representación por impulsos}
	% ============================================================================

	El impulso discreto no necesita ningún aparato distribucional: es una secuencia
	ordinaria, que vale uno en $n=0$ y cero en el resto. Y el escalón vale uno desde
	$n=0$ inclusive:
	\[
		\delta[n]=\begin{cases}1,& n=0\\ 0,& n\neq0\end{cases}
		\qquad\qquad
		u[n]=\begin{cases}1,& n\geq0\\ 0,& n<0\end{cases}
	\]
	Las dos relaciones entre ellos son la primera diferencia y la sumatoria:
	\[
		\delta[n]=u[n]-u[n-1],
		\qquad
		u[n]=\sum_{k=-\infty}^{n}\delta[k].
	\]
	Finalmente, $\delta[n-k]$ vale uno solo en $n=k$, lo que permite escribir cualquier
	secuencia como
	\[
		x[n]=\sum_{k=-\infty}^{\infty}x[k]\,\delta[n-k].
	\]

	% ----------------------------------------------------------------------------
	\ejercicio{11: escalones e impulsos}

	\begin{consigna}
	Sean las secuencias
	\begin{align*}
		x_1[n] &= u[n+2] - u[n-3], \\
		x_2[n] &= u[n+1] + u[n-1] - u[n-3] - u[n-5], \\
		x_3[n] &= n\,\big(u[n] - u[n-5]\big), \\
		x_4[n] &= 2\,\delta[n+1] - \delta[n] + 3\,\delta[n-2].
	\end{align*}
	\begin{enumerate}
		\item Graficar las cuatro secuencias.
		\item Escribir $x_3[n]$ como suma de impulsos desplazados.
	\end{enumerate}
	\end{consigna}

	La técnica para todas las combinaciones de escalones es la misma: identificar desde
	qué índice ``se enciende'' cada escalón y armar la tabla de tramos.

	\begin{enumerate}
		\item \begin{enumerate}[label=\textup{\roman*.}, leftmargin=*, itemsep=4pt]
				\item $x_1[n]=u[n+2]-u[n-3]$. El primer escalón vale uno desde $n=-2$; el
			segundo, desde $n=3$. La resta vale uno donde el primero ya se encendió y el segundo
			todavía no:
			\[
				x_1[n]=1 \ \text{para } -2\leq n\leq 2, \qquad 0 \text{ en el resto}.
			\]
			Es un pulso rectangular discreto de cinco muestras, y su molde general aparece todo
			el tiempo: un pulso de amplitud $A$ que vale $A$ para
			$0\leq n\leq M-1$ y cero fuera se escribe
			\[
				A\big(u[n]-u[n-M]\big).
			\]

			\item $x_2[n]=u[n+1]+u[n-1]-u[n-3]-u[n-5]$. Ahora hay cuatro escalones, que se
			encienden en $n=-1$, $1$, $3$ y $5$. Recorriendo los tramos de izquierda a derecha y
			acumulando:
			\[
				\begin{array}{llll}
					n\leq -2: & 0 & \text{(ninguno encendido)} \\
					n=-1,\,0: & 1 & \text{(solo } u[n+1]) \\
					n=1,\,2:  & 2 & \text{(se suma } u[n-1]) \\
					n=3,\,4:  & 1 & \text{(se resta } u[n-3]) \\
					n\geq 5:  & 0 & \text{(se resta } u[n-5])
				\end{array}
			\]
			El resultado es una escalera que sube un peldaño por vez y luego baja un peldaño por vez,
			simétrica respecto de $n=1{,}5$.

			\item $x_3[n]=n\,(u[n]-u[n-5])$. El paréntesis es el pulso que vale uno para
			$0\leq n\leq 4$, así que el producto deja pasar la rampa $n$ solo en ese tramo:
			\[
				x_3[n] = \{\,\underset{\uparrow}{0}\quad 1\quad 2\quad 3\quad 4\,\}.
			\]
			La muestra en $n=0$ vale cero porque el factor $n$ la anula, no porque el pulso la
			excluya.

			\item $x_4[n]=2\,\delta[n+1]-\delta[n]+3\,\delta[n-2]$. Cada impulso deposita su
			coeficiente en una sola posición: el primero en $n=-1$, el segundo en $n=0$, el tercero
			en $n=2$. Entonces $x_4[-1]=2$, $x_4[0]=-1$, $x_4[2]=3$ y cero en todos los demás
			índices. Las cuatro secuencias están en la Figura~\ref{fig:sol_ej11}.
		\end{enumerate}

		\item La fórmula general
		$x[n]=\sum_k x[k]\,\delta[n-k]$ se aplica escribiendo un término por cada muestra no
		nula, con la muestra como coeficiente y su índice como desplazamiento:
		\[
			\boxed{\;x_3[n] = \delta[n-1] + 2\,\delta[n-2] + 3\,\delta[n-3] + 4\,\delta[n-4].\;}
		\]
		El término de $k=0$ no aparece porque $x_3[0]=0$. Para convencerse de que la igualdad
		vale, basta fijar un índice: en $n=3$ sobrevive un único término ---el de
		$\delta[n-3]$, que vale uno--- y los otros tres se anulan, así que la suma entrega
		$3 = x_3[3]$. A diferencia del tiempo continuo, aquí no hizo falta ningún proceso de
		límite.
	\end{enumerate}

	\begin{figure}[H]
		\centering
		\includestandalone[width=0.88\linewidth]{figures/fig_sol_ej11/fig_sol_ej11}
		\caption{Ejercicio~11. (a) $x_1[n]$, pulso rectangular de cinco muestras.
		(b) $x_2[n]$, escalera. (c) $x_3[n]$, rampa truncada. (d) $x_4[n]$,
		combinación de tres impulsos.}
		\label{fig:sol_ej11}
	\end{figure}

	% ============================================================================
	\tema{El problema del muestreo}
	% ============================================================================

	Muestrear conserva dónde quedó el objeto, no cómo llegó ahí. Si entre dos muestras un
	vector avanza más de media vuelta, las muestras son idénticas a las de otro vector más
	lento que gira en sentido contrario, y no hay manera de distinguirlos. Esta es la
	versión visible del problema; la sección siguiente la convierte en una condición
	exacta sobre el espectro.

	% ----------------------------------------------------------------------------
	\ejercicio{12: la rueda filmada}

	\begin{consigna}
	Una rueda gira a $20$ vueltas por segundo. Sobre el borde lleva
	pintada una marca, que es lo único que permite reconocer en qué
	posición angular está la rueda en cada cuadro. Se la filma a $24$
	cuadros por segundo.
	\begin{enumerate}
		\item Calcular el ángulo que avanza la marca entre cuadro y cuadro.
		\item Reducir ese ángulo al intervalo $(-\pi,\pi]$ y determinar a
		qué velocidad, y en qué sentido, parece girar la rueda en la
		pantalla.
		\item Determinar a qué velocidades de giro la rueda parece quieta.
		\item Determinar la máxima velocidad de giro que esa filmación
		representa sin ambigüedad.
	\end{enumerate}
	\end{consigna}

	Filmar es muestrear: cada cuadro es una foto de la rueda a intervalos regulares de
	\[
		\Delta t = \frac{1}{24}\ \text{s}.
	\]

	\begin{enumerate}
		\item Con $20$ vueltas por segundo, la velocidad angular es
		$\omega = 2\pi\cdot 20 = 40\pi$ rad/s, y entre cuadro y cuadro la marca avanza
		\[
			\lambda = \omega\,\Delta t = \frac{40\pi}{24} = \frac{5\pi}{3}\ \text{rad}
			\;(300^\circ).
		\]
		El mismo número sale sin pasar por radianes: en $1/24$ de segundo una rueda de $20$
		vueltas por segundo completa $20/24 = 5/6$ de vuelta, y cinco sextos de $2\pi$ son
		$5\pi/3$.

		\item La marca avanza casi una vuelta entera por
		cuadro, así que el camino corto entre dos posiciones consecutivas va para el otro lado.
		Reduciendo al intervalo $(-\pi,\pi]$:
		\[
			\frac{5\pi}{3} - 2\pi = -\frac{\pi}{3}\ \text{rad}\;(-60^\circ).
		\]
		El signo negativo dice que la rueda \textbf{parece girar en sentido contrario} al real.
		La velocidad aparente se obtiene volviendo de radianes por cuadro a vueltas por
		segundo: $\pi/3$ rad son $1/6$ de vuelta, y con $24$ cuadros por segundo,
		\[
			\frac{1}{6}\cdot 24 = 4\ \text{vueltas por segundo, hacia atrás}.
		\]
		Una rueda que gira a $20$ vueltas por segundo hacia adelante se ve como una de $4$
		vueltas por segundo hacia atrás.

		\item La rueda aparece quieta si entre cuadro y cuadro la
		marca vuelve exactamente a su posición, es decir si el avance es un número entero de
		vueltas:
		\[
			\lambda = 2\pi m
			\quad\Longleftrightarrow\quad
			\frac{v}{24} = m
			\quad\Longleftrightarrow\quad
			v = 24m \ \text{vueltas por segundo},
		\]
		con $m$ entero. Es decir a $24$, $48$, $72$ vueltas por segundo, y así siguiendo
		---además del caso trivial de la rueda detenida---. La velocidad de filmación queda
		``ciega'' a sus propios múltiplos.

		\item El giro se sigue bien mientras el avance
		por cuadro no llegue a media vuelta:
		\[
			|\lambda| < \pi
			\quad\Longleftrightarrow\quad
			\frac{v}{24} < \frac{1}{2}
			\quad\Longleftrightarrow\quad
			v < 12 \ \text{vueltas por segundo}.
		\]
		El tope es la mitad de la velocidad de filmación, $12$ vueltas por segundo, y no se
		alcanza: justo ahí la marca da media vuelta por cuadro y no se puede decidir si giró
		hacia un lado o hacia el otro. Ese tope es la frecuencia de Nyquist de esta filmación.
	\end{enumerate}

	% ============================================================================
	\tema{El espectro de una señal muestreada}
	% ============================================================================

	Modelando el muestreo como el producto de la señal por un tren de impulsos de período
	$\Delta t$, el espectro de la señal muestreada resulta
	\[
		X_s(j\omega) = \frac{1}{\Delta t}\sum_{k=-\infty}^{\infty}
		               X\!\big(j(\omega-k\omega_s)\big),
		\qquad \omega_s=\frac{2\pi}{\Delta t}.
	\]
	En palabras: \textbf{muestrear replica el espectro}. Aparece una copia de $X(j\omega)$
	centrada en cada múltiplo de la frecuencia de muestreo, todas con la misma forma y
	todas escaladas por el mismo factor $1/\Delta t$. Como consecuencia inmediata,
	$X_s(j\omega)$ es siempre periódico de período $\omega_s$.

	% ----------------------------------------------------------------------------
	\ejercicio{13: replicación de un espectro triangular}

	\begin{consigna}
	Una señal $x(t)$ tiene espectro $X(j\omega)$ triangular, de altura
	$1$ en el origen y nulo fuera de $|\omega| < 3000\pi$ rad/s.
	\begin{enumerate}
		\item Con $\Delta t = 0{,}2$ ms, calcular $\omega_s$ y graficar
		$X_s(j\omega)$ mostrando al menos tres copias, con sus posiciones
		y su factor de escala. Demarcar la banda base y verificar sobre la
		gráfica que $X_s(j\omega)$ es periódico de período $\omega_s$.
		\item Repetir con $\Delta t = 0{,}4$ ms e indicar en qué banda de
		frecuencias se produce el solapamiento.
		\item Repetir las dos gráficas con el eje medido en frecuencia
		digital $\lambda = \omega\,\Delta t$, indicando en cada caso cuánto
		ocupa una copia y dónde queda la banda base. Escribir la condición
		de Nyquist en términos de $\lambda$.
	\end{enumerate}
	\end{consigna}

	La señal está limitada en banda con $\omega_M = 3000\pi$ rad/s.

	\begin{enumerate}
		\item La frecuencia de muestreo es
		\[
			\omega_s = \frac{2\pi}{\Delta t} = \frac{2\pi}{0{,}2\times10^{-3}}
			         = 10\,000\pi \ \text{rad/s} \approx 31\,416\ \text{rad/s},
		\]
		y el factor de escala, $1/\Delta t = 5000$. El espectro muestreado es entonces una
		sucesión de triángulos idénticos al original, de altura $5000$ en su centro y base
		$\pm 3000\pi$ alrededor de él, centrados en
		\[
			\omega = 0,\ \pm 10\,000\pi,\ \pm 20\,000\pi,\ \dots
		\]
		El panel (a) de la Figura~\ref{fig:sol_ej13} muestra cinco de ellos. Las copias no se
		tocan, y se puede medir cuánto sobra: entre el borde derecho de la copia central
		($3000\pi$) y el borde izquierdo de la vecina ($10\,000\pi - 3000\pi = 7000\pi$)
		quedan $4000\pi$ rad/s de separación. Es lo que asegura la condición de Nyquist,
		$\omega_s = 10\,000\pi > 2\omega_M = 6000\pi$.

		\emph{La banda base.} Es la franja $|\omega| < \omega_s/2 = 5000\pi$ rad/s, sombreada en
		gris en la figura. La determina la frecuencia de muestreo sola, sin que la señal
		intervenga, y es el tramo del eje que la secuencia efectivamente guarda: lo que hay
		afuera es copia de lo que hay adentro. La copia central llega hasta $3000\pi$, así que
		entra completa y todavía sobran $2000\pi$ rad/s hasta el borde.

		\emph{Periodicidad.} Sobre la gráfica se verifica de inmediato: todas las copias son
		iguales y sus centros están espaciados exactamente $\omega_s$, de modo que desplazar el
		eje en $\omega_s$ deja cada copia sobre la posición de su vecina y la figura queda como
		estaba. En símbolos, $X_s(j(\omega+\omega_s)) = X_s(j\omega)$.

		\item Al duplicar el intervalo, la
		frecuencia de muestreo se reduce a la mitad y el factor de escala también:
		\[
			\omega_s = \frac{2\pi}{0{,}4\times10^{-3}} = 5000\pi\ \text{rad/s}
			\approx 15\,708\ \text{rad/s},
			\qquad \frac{1}{\Delta t} = 2500 .
		\]
		Ahora $2\omega_M = 6000\pi > 5000\pi = \omega_s$: la condición de Nyquist no se cumple
		y las copias se superponen. Para ubicar dónde, alcanza con comparar los bordes de dos
		copias vecinas. La central llega hasta $3000\pi$; la centrada en $\omega_s$ empieza en
		\[
			\omega_s - \omega_M = 5000\pi - 3000\pi = 2000\pi .
		\]
		Como $2000\pi < 3000\pi$, la vecina empieza antes de que termine la central, y las dos
		comparten el tramo
		\[
			\boxed{\;2000\pi < |\omega| < 3000\pi \quad (\text{aproximadamente } 6283 <
			|\omega| < 9425 \ \text{rad/s}).\;}
		\]
		Ahí los dos espectros se suman, y el espectro muestreado queda por encima del
		original: el valle que en el caso anterior bajaba hasta cero ahora aparece relleno
		con contenido que no pertenece ahí. Es lo que se ve en el panel (b).

		\item Cambiar de unidad es
		multiplicar cada frecuencia por $\Delta t$. Lo primero que cambia es la separación entre
		copias, que en los dos casos pasa a ser la misma:
		\[
			\omega_s\,\Delta t = \frac{2\pi}{\Delta t}\,\Delta t = 2\pi .
		\]
		Con esa separación queda fijada también la banda base, que en frecuencia digital es
		siempre $[-\pi,\pi]$, cualquiera sea el intervalo de muestreo.

		Lo que sí distingue a los dos casos es cuánto ocupa una copia. El borde de la copia
		central, $\omega_M = 3000\pi$ rad/s, cae en
		\[
			\lambda_M = \omega_M\,\Delta t =
			\begin{cases}
				3000\pi\cdot 0{,}2\times10^{-3} = 0{,}6\pi, & \Delta t = 0{,}2\ \text{ms},\\[4pt]
				3000\pi\cdot 0{,}4\times10^{-3} = 1{,}2\pi, & \Delta t = 0{,}4\ \text{ms}.
			\end{cases}
		\]
		La Figura~\ref{fig:sol_ej13_lambda} muestra las dos gráficas. Con $0{,}6\pi$ la copia
		entra en la banda base y sobran $0{,}4\pi$ hasta el borde; con $1{,}2\pi$ se pasa de
		largo y se mete en el territorio de la vecina, que empieza en $2\pi - 1{,}2\pi =
		0{,}8\pi$. El tramo compartido va de $0{,}8\pi$ a $1{,}2\pi$, y es el mismo solapamiento
		del inciso anterior medido en la unidad nueva, porque $2000\pi\cdot0{,}4$ ms $=0{,}8\pi$
		y $3000\pi\cdot0{,}4$ ms $=1{,}2\pi$.

		Leída así, la condición de Nyquist se escribe en una línea,
		\[
			\boxed{\;\lambda_M < \pi\;}
		\]
		que en palabras es que la copia entre en la banda base. Es la condición de siempre en
		otra unidad, porque $\lambda_M = \omega_M\,\Delta t < \pi$ equivale a
		$\omega_M < \pi/\Delta t = \omega_s/2$.
	\end{enumerate}

	\begin{figure}[H]
		\centering
		\includestandalone[width=0.92\linewidth]{figures/fig_sol_ej13/fig_sol_ej13}
		\caption{Ejercicio~13. Espectro muestreado, en los dos casos, sobre el mismo
		eje de frecuencias. La franja gris con la flecha es la banda base,
		$|\omega|<\omega_s/2$. (a) $\Delta t = 0{,}2$ ms: las copias, de altura $5000$ y
		espaciadas $\omega_s = 10\,000\pi$, quedan separadas. (b) $\Delta t = 0{,}4$
		ms: las copias, de altura $2500$ y espaciadas $\omega_s = 5000\pi$, se
		solapan; en las zonas sombreadas los flancos de dos copias vecinas se suman
		(punteado gris).}
		\label{fig:sol_ej13}
	\end{figure}

	\begin{figure}[H]
		\centering
		\includestandalone[width=0.92\linewidth]{figures/fig_sol_ej13_lambda/fig_sol_ej13_lambda}
		\caption{Ejercicio~13, inciso (c). Los dos casos con el eje medido en frecuencia
		digital. Las copias quedan separadas $2\pi$ en los dos, y la banda base ---la
		franja gris--- es siempre $[-\pi,\pi]$. Lo único que cambia es cuánto ocupa una
		copia. (a) $\Delta t = 0{,}2$ ms: la copia llega a $0{,}6\pi$ y entra holgada.
		(b) $\Delta t = 0{,}4$ ms: llega a $1{,}2\pi$, se pasa del borde y se solapa con la
		vecina sobre las zonas sombreadas.}
		\label{fig:sol_ej13_lambda}
	\end{figure}

	% ============================================================================
	\tema{Nyquist y aliasing}
	% ============================================================================

	Si $x(t)$ está limitada en banda a $\omega_M$, las copias del espectro no se solapan
	cuando
	\[
		\omega_s > 2\,\omega_M .
	\]
	Las dos cantidades tienen nombre propio y no conviene confundirlas:
	\begin{itemize}[nosep]
		\item \textbf{Velocidad de Nyquist} $= 2\omega_M$: la frecuencia de muestreo mínima que
		      hay que superar. Es una propiedad \emph{de la señal}.
		\item \textbf{Frecuencia de Nyquist} $= \omega_s/2$: la frecuencia más alta que ese
		      muestreo puede representar, es decir, el borde de la banda base. Es una propiedad
		      \emph{del muestreo}.
	\end{itemize}
	Cuando la condición se viola, una componente de frecuencia $\omega_0$ por encima de
	$\omega_s/2$ reaparece dentro de la banda base en su \textbf{frecuencia alias}
	\[
		\omega_a = |\,\omega_0 - k\,\omega_s\,|,
	\]
	con el entero $k$ que más la acerca al origen.

	% ----------------------------------------------------------------------------
	\ejercicio{14: frecuencia máxima y velocidad de Nyquist}

	\begin{consigna}
	Para cada una de las siguientes señales, determinar la frecuencia
	máxima $\omega_M$ y la velocidad de Nyquist, y proponer una frecuencia de
	muestreo con margen sobre esa velocidad, indicando el intervalo $\Delta t$
	correspondiente.
	\begin{enumerate}
		\item $x(t) = 10\cos(5t) - 5\sin(7\pi t) + 3\cos(30t)$
		\item $y(t) = 18\sin(10t) + 15\sin(5\pi t) + 5\cos(60t)$
		\item $z(t) = 15\cos(6\pi t) + 7\sin(10\pi t) + 2\cos(15\pi t)$
	\end{enumerate}
	\end{consigna}

	En una suma de sinusoides, la frecuencia máxima es simplemente la mayor de las
	frecuencias presentes. El único cuidado es leerlas bien: lo que está multiplicando a
	$t$ dentro del coseno o del seno ya es $\omega$ en rad/s, y las amplitudes no
	intervienen.

	\begin{enumerate}
		\item $x(t)=10\cos(5t)-5\sin(7\pi t)+3\cos(30t)$. Las frecuencias son $5$,
		$7\pi \approx 21{,}99$ y $30$ rad/s. La mayor es
		\[
			\omega_M = 30 \ \text{rad/s}
			\quad\Longrightarrow\quad
			2\omega_M = 60 \ \text{rad/s}.
		\]
		Adoptando un margen del $50\%$, $\omega_s = 90$ rad/s y
		$\Delta t = 2\pi/90 \approx 69{,}8$ ms.

		\item $y(t)=18\sin(10t)+15\sin(5\pi t)+5\cos(60t)$. Frecuencias $10$,
		$5\pi\approx 15{,}71$ y $60$ rad/s, de modo que
		\[
			\omega_M = 60\ \text{rad/s}
			\quad\Longrightarrow\quad
			2\omega_M = 120\ \text{rad/s},
		\]
		y con el mismo criterio $\omega_s = 180$ rad/s, $\Delta t \approx 34{,}9$ ms.

		\item $z(t)=15\cos(6\pi t)+7\sin(10\pi t)+2\cos(15\pi t)$. Frecuencias $6\pi$,
		$10\pi$ y $15\pi$ rad/s. La mayor es
		\[
			\omega_M = 15\pi \approx 47{,}12\ \text{rad/s}
			\quad\Longrightarrow\quad
			2\omega_M = 30\pi \approx 94{,}25 \ \text{rad/s}.
		\]
		Adoptando $\omega_s = 45\pi \approx 141{,}4$ rad/s queda
		$\Delta t = 2\pi/45\pi = 2/45 \approx 44{,}4$ ms. Notar que la componente más rápida es
		la de amplitud más pequeña: la amplitud no influye en la velocidad de Nyquist.
	\end{enumerate}

	% ----------------------------------------------------------------------------
	\ejercicio{15: aliasing de una componente}

	\begin{consigna}
	Para la señal $z(t)$ del ejercicio anterior se adopta
	$\omega_s = 80$ rad/s.
	\begin{enumerate}
		\item Verificar si se cumple la condición de Nyquist.
		\item Indicar qué componentes quedan bien muestreadas.
		\item Calcular la frecuencia alias de la componente mal muestreada
		e indicar dónde queda respecto de las componentes bien muestreadas.
		\item Indicar si esa componente puede recuperarse con un
		procesamiento posterior. Justificar.
		\item Determinar la mínima frecuencia de muestreo que evita el
		problema.
	\end{enumerate}
	\end{consigna}

	La señal es la del inciso (c) anterior, con componentes en $6\pi$, $10\pi$ y $15\pi$
	rad/s, y se muestrea con $\omega_s = 80$ rad/s.

	\begin{enumerate}
		\item La velocidad de Nyquist de esta señal es
		$2\omega_M = 30\pi \approx 94{,}25$ rad/s, y
		\[
			\omega_s = 80 \;<\; 94{,}25 = 2\omega_M .
		\]
		\textbf{No se cumple}: se está muestreando por debajo de lo necesario.

		\item El criterio, componente por componente, es
		caer por debajo de la frecuencia de Nyquist $\omega_s/2 = 40$ rad/s:
		\[
			6\pi \approx 18{,}85 < 40 \ \checkmark,
			\qquad
			10\pi \approx 31{,}42 < 40\ \checkmark,
			\qquad
			15\pi \approx 47{,}12 > 40 \ \times .
		\]
		Las dos primeras quedan bien muestreadas; la de $15\pi$ es la única que viola la
		condición, y es la que causa todo el problema.

		\item La componente de $\omega_0 = 15\pi$ genera réplicas
		en $\omega_0 - k\omega_s$; hay que elegir el $k$ que la acerque más al origen, que
		aquí es $k=1$:
		\[
			\omega_a = |\,15\pi - 80\,| = |47{,}12 - 80| = 32{,}88 \ \text{rad/s}.
		\]
		Después de muestrear, la componente de $47{,}12$ rad/s se presenta como si fuera una de
		$32{,}88$ rad/s.

		Dónde cae es lo preocupante. El alias queda a $1{,}46$ rad/s de $10\pi \approx 31{,}42$
		rad/s, que es una componente legítima de la señal, mientras que la otra componente bien
		muestreada, $6\pi \approx 18{,}85$ rad/s, queda lejos. El intruso se instala
		prácticamente encima de un vecino real.

		\item \textbf{No.} Después del muestreo, el alias
		ocupa \emph{exactamente} la frecuencia $32{,}88$ rad/s, indistinguible de una
		componente genuina de esa frecuencia que la señal podría haber tenido. Ningún
		procesamiento posterior ---analógico o digital--- tiene información para separarlas,
		porque en las muestras las dos son indistinguibles. El daño es irreversible, y por eso la
		única defensa es actuar antes: o se filtra la componente rápida en tiempo continuo, o
		se aumenta la frecuencia de muestreo.

		\item La condición es
		\[
			\omega_s > 2\omega_M = 30\pi \approx 94{,}25\ \text{rad/s},
		\]
		que en la práctica se satisface con algo de margen, por ejemplo $\omega_s = 120$ rad/s.
	\end{enumerate}

	% ----------------------------------------------------------------------------
	\ejercicio{16: solapamiento de un espectro trapezoidal}

	\begin{consigna}
	El espectro de amplitud de una señal $x(t)$ es el de la
	Figura~\ref{fig:tp6_espectro_trapecio}, nulo fuera de $|\omega| < 500$
	rad/s. Se la muestrea con una frecuencia de muestreo $f_s = 150$ Hz.

	\begin{center}
		\includestandalone[width=0.72\linewidth]{../figures/fig_tp6_espectro_trapecio/fig_tp6_espectro_trapecio}
		\captionof{figure}{Espectro de amplitud del ejercicio 16.}
		\label{fig:tp6_espectro_trapecio}
	\end{center}

	\begin{enumerate}
		\item Calcular $\omega_s$ y la frecuencia de Nyquist, demarcar la
		banda base y verificar si se cumple la condición de no
		solapamiento.
		\item Determinar qué porción del espectro se solapa y sobre qué
		banda de la banda base se deposita.
		\item Graficar el espectro muestreado sobre la banda base,
		señalando la zona contaminada.
		\item Analizar cómo se ve la componente de $\omega = 460$ rad/s
		después del muestreo.
		\item Determinar la mínima $f_s$ que evita el solapamiento.
	\end{enumerate}
	\end{consigna}

	El espectro es nulo fuera de $|\omega| < 500$ rad/s, así que $\omega_M = 500$ rad/s.

	\begin{enumerate}
		\item Con $f_s = 150$ Hz,
		\[
			\omega_s = 2\pi f_s = 300\pi \approx 942{,}5 \ \text{rad/s},
			\qquad
			\frac{\omega_s}{2} = 150\pi \approx 471{,}2\ \text{rad/s},
		\]
		de modo que la banda base es la franja $|\omega| < 471{,}2$ rad/s, que es lo que la
		secuencia guarda. La velocidad de Nyquist de la señal es $2\omega_M = 1000$ rad/s, y
		\[
			\omega_s = 942{,}5 \;<\; 1000 = 2\omega_M .
		\]
		La condición de no solapamiento \textbf{no se cumple}, aunque por poco: bastaría con
		aumentar la frecuencia de muestreo algo más del $6\%$.

		\item Hay dos preguntas distintas y conviene
		separarlas.

		\emph{Dónde se solapan las copias.} La copia central ocupa hasta $500$ rad/s y la
		vecina, centrada en $\omega_s$, empieza en
		\[
			\omega_s - \omega_M = 942{,}5 - 500 = 442{,}5 \ \text{rad/s},
		\]
		de modo que las dos comparten el tramo $442{,}5 < \omega < 500$ rad/s.

		\emph{Qué porción del original se pliega sobre la banda base.} Es la que está por
		encima de la frecuencia de Nyquist: el contenido entre $471{,}2$ y $500$ rad/s. Cada
		componente de frecuencia $\omega_0$ en ese tramo reaparece en
		$\omega_a = \omega_s - \omega_0$, y a medida que $\omega_0$ recorre $(471{,}2;\,500)$,
		su alias recorre $(442{,}5;\,471{,}2)$ ---en orden invertido, porque el pliegue actúa
		como un espejo ubicado en $\omega_s/2$---. La zona contaminada de la banda base es
		entonces
		\[
			\boxed{\;442{,}5 < |\omega| < 471{,}2 \ \text{rad/s}.\;}
		\]

		\item La Figura~\ref{fig:sol_ej16} muestra las dos vistas: el
		espectro muestreado completo con sus copias solapadas, y el detalle de la banda base
		con la zona contaminada.

		\item Como $460 < 471{,}2$, esta
		componente está por debajo de la frecuencia de Nyquist: su propia réplica cae donde
		corresponde y no se convierte en alias de nada. Pero $460$ está dentro del tramo
		contaminado, así que \emph{recibe} el alias de otra componente. Cuál, se despeja de la
		relación del pliegue:
		\[
			\omega_0 = \omega_s - \omega_a = 942{,}5 - 460 = 482{,}5\ \text{rad/s}.
		\]
		Poniendo números sobre el trapecio, cuyo flanco baja linealmente de $1$ en $300$ rad/s
		a $0$ en $500$ rad/s ---es decir, $|X| = (500-\omega)/200$---,
		\[
			\text{valor verdadero en } 460: \ \frac{500-460}{200}=0{,}200,
			\qquad
			\text{alias que se le suma}: \ \frac{500-482{,}5}{200}=0{,}0875 .
		\]
		Después de muestrear se mide $0{,}200+0{,}0875 = 0{,}2875$. La componente sobrevive,
		pero con un error de amplitud del $44\%$, y no hay forma de separar las dos
		contribuciones.

		\item Se pide $\omega_s > 2\omega_M = 1000$
		rad/s, es decir
		\[
			f_s > \frac{1000}{2\pi} = 159{,}2\ \text{Hz}.
		\]
		Basta aumentar levemente la frecuencia de muestreo ---por ejemplo a $200$ Hz--- para que
		el problema desaparezca por completo.
	\end{enumerate}

	\begin{figure}[H]
		\centering
		\includestandalone[width=0.92\linewidth]{figures/fig_sol_ej16/fig_sol_ej16}
		\caption{Ejercicio~16. (a) Espectro muestreado (azul) como suma de las copias
		(punteado gris) espaciadas $\omega_s = 942{,}5$ rad/s; en las zonas sombreadas
		las copias se solapan. (b) Detalle de la banda base: en punteado, el espectro
		original; en azul, lo que el muestreo entrega. Difieren solo en la zona
		sombreada, donde el espectro medido queda por encima del verdadero. La
		componente de $460$ rad/s debería valer $0{,}200$ y se mide $0{,}2875$. Las
		alturas están normalizadas por $1/\Delta t$.}
		\label{fig:sol_ej16}
	\end{figure}

	% ============================================================================
	\tema{Del modelo a la práctica}
	% ============================================================================

	La cadena real incorpora a la entrada un bloque que el modelo ideal no menciona, el
	\textbf{filtro antialiasing}: un pasabajos analógico que recorta lo que supere la
	frecuencia de Nyquist \emph{antes} de muestrear.

	% ----------------------------------------------------------------------------
	\ejercicio{17: señal útil con ruido fuera de banda}

	\begin{consigna}
	Una señal útil ocupa la banda $|\omega| < \omega_{\max}$ y se le
	suma un ruido concentrado en $\omega_1 < |\omega| < \omega_2$, según la
	Figura~\ref{fig:tp6_senal_ruido}. Tomar
	$\omega_{\max} = 1000$ rad/s, $\omega_1 = 1500$ rad/s y
	$\omega_2 = 2000$ rad/s.

	\begin{center}
		\includestandalone[width=0.72\linewidth]{../figures/fig_tp6_senal_ruido/fig_tp6_senal_ruido}
		\captionof{figure}{Espectro de la señal útil (A) y del ruido (B),
		ejercicio 17.}
		\label{fig:tp6_senal_ruido}
	\end{center}

	\begin{enumerate}
		\item Con $\omega_s = 2500$ rad/s, demarcar la banda base, ubicar
		las réplicas del ruido y determinar si contaminan la banda útil.
		\item Proponer dos maneras distintas de muestrear sin afectar la
		banda útil, y para cada una indicar qué condición debe cumplir
		$\omega_s$.
		\item Explicar por qué el filtro antialiasing debe ser analógico y
		actuar antes del muestreo.
		\item Explicar por qué en la práctica la frecuencia de corte se
		ubica algo por debajo de la de Nyquist.
	\end{enumerate}
	\end{consigna}

	La señal útil ocupa $|\omega| < 1000$ rad/s y el ruido, $1500 < |\omega| < 2000$ rad/s.
	Notar que el ruido está \emph{separado} de la señal en frecuencia: es lo que hace
	posible el ejercicio.

	\begin{enumerate}
		\item \emph{La banda base.} Es la franja $|\omega| < \omega_s/2 = 1250$ rad/s, que es lo que
		la secuencia guarda. La banda útil llega hasta $\omega_{\max} = 1000$ rad/s, así que
		entra entera y le sobra una franja libre de $250$ rad/s hasta el borde.

		\emph{Las réplicas del ruido.} El muestreo repite el espectro completo ---señal útil y
		ruido juntos--- centrado en cada múltiplo de $\omega_s$. Ese espectro completo se
		extiende hasta
		\[
			\omega_M = \omega_2 = 2000\ \text{rad/s},
		\]
		de modo que su velocidad de Nyquist es $2\omega_M = 4000$ rad/s y $\omega_s = 2500$ no
		la alcanza. La condición de no solapamiento no se cumple, y hay que ubicar las réplicas
		para ver qué se pisa con qué.

		Cada una de las dos bandas de ruido reaparece desplazada un múltiplo de $\omega_s$. La
		Tabla~\ref{tab:sol_ej17} las lista para los desplazamientos que dejan algo cerca de la
		banda base.

		\begin{table}[H]
			\centering
			\caption{Réplicas de las dos bandas de ruido, desplazadas $k\,\omega_s$ con
			$\omega_s = 2500$ rad/s. En negrita, las que caen dentro de la banda base
			$|\omega| < 1250$ rad/s.}
			\label{tab:sol_ej17}
			\begin{tabular}{ccc}
				\toprule
				$k$ & $[1500,\,2000] + k\,\omega_s$ & $[-2000,\,-1500] + k\,\omega_s$ \\
				\midrule
				$-1$ & $\mathbf{[-1000,\,-500]}$ & $[-4500,\,-4000]$ \\
				$0$  & $[1500,\,2000]$           & $[-2000,\,-1500]$ \\
				$+1$ & $[4000,\,4500]$           & $\mathbf{[500,\,1000]}$ \\
				\bottomrule
			\end{tabular}
		\end{table}

		De todas ellas, las únicas que caen dentro de la banda base son $[500,\,1000]$ y
		$[-1000,\,-500]$, que vienen de desplazar la banda de ruido opuesta en $\pm\omega_s$.
		La Figura~\ref{fig:sol_ej17} las muestra en punteado, junto con las copias vecinas
		completas.

		\emph{Si contaminan la banda útil.} La banda útil ocupa $|\omega| < 1000$ rad/s, y esas
		dos réplicas quedan enteramente adentro, superpuestas al flanco de la banda útil:
		\[
			\boxed{\;\text{sí, el ruido contamina la banda útil, sobre }
			500 < |\omega| < 1000 \ \text{rad/s}.\;}
		\]
		Sobre ese tramo la señal muestreada es la suma del flanco de la banda útil y la réplica
		del ruido, y una vez sumadas ocupan las mismas frecuencias.

		La señal útil, por sí sola, estaría bien muestreada, porque su velocidad de Nyquist es
		$2\omega_{\max} = 2000$ rad/s y $\omega_s = 2500$ la supera. Lo que rompe la condición
		es el ruido, que empuja la frecuencia máxima de $1000$ a $2000$ rad/s y con ella la
		velocidad de Nyquist de $2000$ a $4000$ rad/s.

		\item Las dos estrategias atacan el problema desde
		lados opuestos: una elimina el ruido, la otra le hace lugar.

		\emph{Primera: filtrar antes de muestrear.} Se intercala un pasabajos analógico
		---el filtro antialiasing--- con frecuencia de corte entre $\omega_{\max}=1000$ y
		$\omega_1=1500$ rad/s, que deja pasar la banda útil entera y elimina el ruido. Con el
		ruido fuera, lo único que se muestrea es la señal útil, y basta pedir
		\[
			\omega_s > 2\,\omega_{\max} = 2000 \ \text{rad/s}.
		\]
		Es la solución económica: permite el muestreo más lento posible.

		\emph{Segunda: muestrear a frecuencia suficientemente alta.} Sin filtrar nada, se elige $\omega_s$ de
		modo que ninguna réplica del ruido alcance la banda útil. La réplica más peligrosa es
		la de la banda positiva desplazada en $-\omega_s$, que ocupa
		$[\omega_1-\omega_s,\ \omega_2-\omega_s]$; su borde derecho tiene que quedar por debajo
		de $-\omega_{\max}$:
		\[
			\omega_2 - \omega_s \leq -\omega_{\max}
			\quad\Longleftrightarrow\quad
			\omega_s \geq \omega_2 + \omega_{\max} = 2000+1000 = 3000\ \text{rad/s}.
		\]
		La otra banda produce la misma condición por simetría. Con $\omega_s \geq 3000$ rad/s
		el ruido sigue estando ---cae entre $1000$ rad/s y la frecuencia de Nyquist---, pero ya
		no se mezcla con lo útil, y un filtro digital posterior puede quitarlo. El costo es
		muestrear un $50\%$ más rápido y procesar más datos.

		Las dos alternativas conviven solo porque el enunciado plantea un ruido acotado,
		confinado entre $\omega_1$ y $\omega_2$. Una señal que proviene de un
		fenómeno físico no ofrece esa garantía. El ruido térmico de la electrónica, la
		interferencia de la red y lo que captan los cables no se terminan en ninguna frecuencia
		en particular, de modo que por encima de cualquier $\omega_s/2$ que se elija siempre
		queda algo de contenido. Con un espectro que no termina, la segunda estrategia se queda
		sin condición que cumplir, porque ninguna $\omega_s$ deja todas las réplicas fuera de
		la banda útil, y lo que quede más allá de la frecuencia de Nyquist se pliega igual sobre
		la banda base.

		Por eso, en una cadena de adquisición concreta el filtro antialiasing es obligatorio,
		porque es lo único que acota el espectro que llega al muestreador. La elección de $\omega_s$ pasa a ser complementaria. Aumentarla le deja al
		filtro una banda de transición más ancha, y con eso el filtro puede ser de menor orden,
		pero no lo reemplaza.

		\item Porque después del
		muestreo ya es tarde. Una vez que las réplicas del ruido se sumaron sobre la banda
		útil, la señal y el alias ocupan exactamente las mismas frecuencias, y no queda ninguna
		característica que permita separarlos: cualquier filtro que borre el alias borra
		también la parte de la señal que está ahí. La única ventana de oportunidad es mientras
		el ruido todavía está en su frecuencia original, separado de lo útil, y eso ocurre solo
		en tiempo continuo, antes del conversor. En este ejercicio se ve con claridad: antes de
		muestrear, el ruido vive en $[1500,2000]$ y basta un pasabajos común para sacarlo;
		después de muestrear, vive en $[500,1000]$, encima de la señal.

		\item Porque
		los filtros reales no cortan de manera abrupta: entre la banda de paso y la banda de
		rechazo hay una \emph{banda de transición} donde la atenuación crece gradualmente. Si el
		corte se pusiera justo en $\omega_s/2$, todo lo que cae en esa transición pasaría
		atenuado pero no nulo, y al muestrear se plegaría igual sobre la banda base. Dejando el
		corte algo más abajo, la transición completa queda por debajo de $\omega_s/2$ y llega
		al muestreador ya suficientemente atenuada. El precio es resignar algo de ancho de
		banda útil, o equivalentemente muestrear a una frecuencia algo mayor, a cambio de evitar
		un daño irreparable.
	\end{enumerate}

	\begin{figure}[H]
		\centering
		\includestandalone[width=0.95\linewidth]{figures/fig_sol_ej17/fig_sol_ej17}
		\caption{Ejercicio~17, inciso (a). Espectro muestreado con
		$\omega_s = 2500$ rad/s. En trazo lleno, la copia central ---señal útil en azul y
		ruido en rojo---; en punteado, las copias vecinas. Como $\omega_s$ no llega a
		$2\omega_M = 4000$ rad/s, la copia central y la vecina comparten la franja marcada
		abajo, y las réplicas del ruido caen sobre $500 < |\omega| < 1000$ (zonas
		sombreadas), dentro de la banda útil. La señal por sí sola cumple Nyquist
		($\omega_s/2 = 1250 > 1000$).}
		\label{fig:sol_ej17}
	\end{figure}

	% ============================================================================
	\tema{El eje de frecuencias medido en $\lambda$}
	% ============================================================================

	Medir la frecuencia en $\lambda = \omega\,\Delta t$ es cambiar la unidad del eje
	horizontal, de radianes por segundo a radianes por muestra. El cambio reescala el eje, y
	un paso de $\omega_s$ pasa a ser un paso de
	\[
		\omega_s\,\Delta t = 2\pi .
	\]
	Sobre el eje $\lambda$, entonces, las copias del espectro quedan separadas $2\pi$, la
	banda base ocupa $[-\pi,\pi]$ y la frecuencia de Nyquist $\omega_s/2$ cae en
	$\lambda = \pi$, sea cual sea el intervalo de muestreo. Para pasar de una unidad a la
	otra:
	\[
		\lambda = \omega\,\Delta t = \frac{2\pi f}{f_s},
		\qquad\qquad
		f = \frac{\lambda\, f_s}{2\pi}.
	\]

	% ----------------------------------------------------------------------------
	\ejercicio{18: dos señales sobre el mismo eje $\lambda$}

	\begin{consigna}
	Una señal de audio abarca de $20$ Hz a $20$ kHz y se muestrea
	con $f_s = 44{,}1$ kHz. Una señal de vibración mecánica abarca de
	$0{,}5$ Hz a $500$ Hz.
	\begin{enumerate}
		\item Ubicar los extremos de la banda de audio sobre el eje
		$\lambda$ y demarcar la banda base.
		\item Determinar la frecuencia de muestreo que hace que la señal
		de vibración ocupe exactamente el mismo tramo del eje $\lambda$,
		y el intervalo $\Delta t$ correspondiente.
		\item Representar sobre el eje $\lambda$ la banda que ocupa cada
		una de las dos señales muestreadas y comparar las dos
		representaciones.
		\item Calcular, en Hz, la frecuencia continua que corresponde a
		$\lambda = \pi/3$ y a $\lambda = \pi$ en cada uno de los dos
		casos, y explicar por qué el mismo $\lambda$ representa
		frecuencias tan distintas.
	\end{enumerate}
	\end{consigna}

	\begin{enumerate}
		\item Con $f_s = 44{,}1$ kHz se
		aplica $\lambda = 2\pi f/f_s$ a los dos extremos:
		\[
			\lambda_{\min} = \frac{2\pi\cdot 20}{44\,100} = 2{,}85\times10^{-3}\ \text{rad}
			\ (0{,}16^\circ),
			\qquad
			\lambda_{\max} = \frac{2\pi\cdot 20\,000}{44\,100} = 2{,}849\ \text{rad}
			\ (163{,}3^\circ).
		\]
		Expresado en fracciones de $\pi$, la banda va de $0{,}0009\pi$ a $0{,}907\pi$: empieza
		casi en el origen y llega casi hasta $\pi$. La banda base es el tramo $[-\pi,\pi]$, y
		el panel (a) de la Figura~\ref{fig:sol_ej18} la muestra sombreada, con la señal adentro
		y las imágenes que el muestreo deja a los costados, separadas $2\pi$.

		Que la banda llegue tan cerca del borde tiene una explicación conocida. La velocidad de
		Nyquist del audio es $2\cdot 20 = 40$ kHz y la frecuencia de muestreo es $44{,}1$ kHz,
		apenas un $10{,}25\%$ por encima; ese margen chico en hertz es el mismo margen chico
		que se ve entre $0{,}907\pi$ y $\pi$.

		\item Se pide que la señal de vibración
		llegue al mismo $\lambda_{\max}$, así que se iguala y se despeja la frecuencia de
		muestreo:
		\[
			\frac{2\pi\cdot 500}{f_s'} = \frac{2\pi\cdot 20\,000}{44\,100}
			\quad\Longrightarrow\quad
			\boxed{\;f_s' = 500\cdot\frac{44\,100}{20\,000} = 1102{,}5\ \text{Hz},
			\qquad \Delta t' = \frac{1}{1102{,}5} = 0{,}907\ \text{ms}.\;}
		\]
		El extremo inferior se acomoda solo:
		\[
			\frac{2\pi\cdot 0{,}5}{1102{,}5} = 2{,}85\times10^{-3}\ \text{rad},
		\]
		el mismo valor que en el audio. Y tiene que ser así, porque las dos bandas están en la
		misma proporción ---$500/0{,}5 = 20\,000/20 = 1000$--- y las dos frecuencias de
		muestreo también: $1102{,}5$ es $44\,100$ dividido $40$, el mismo factor que separa
		$500$ Hz de $20$ kHz. Al dividir por $40$ las frecuencias y la frecuencia de muestreo,
		el cociente $f/f_s$ de cada extremo no cambia, y $\lambda$ depende solo de ese
		cociente. El intervalo de muestreo, en cambio, se estira cuarenta veces: $22{,}7\
		\mu$s para el audio contra $0{,}907$ ms para la vibración.

		\item Son la misma, y ese es el punto. Los dos
		paneles de la Figura~\ref{fig:sol_ej18} tienen el dibujo idéntico: la misma banda
		ocupada entre $-0{,}907\pi$ y $0{,}907\pi$, la misma banda base, las mismas imágenes
		separadas $2\pi$. Lo único que los distingue es la regla en hertz que va debajo del
		eje, donde $\lambda = \pi$ vale $22{,}05$ kHz en un caso y $551{,}25$ Hz en el otro.

		Dicho de otro modo, las dos secuencias son indistinguibles mientras se las mire como
		secuencias. Nada en ellas registra que una viene de un micrófono y la otra de un
		acelerómetro, ni que los fenómenos que las produjeron transcurren a ritmos que
		difieren en un factor de cuarenta.

		\item Se despeja $f$ de
		$\lambda = 2\pi f/f_s$, con $f = \lambda f_s/2\pi$:
		\[
			\lambda=\frac{\pi}{3}
			\ \Longrightarrow\
			f = \frac{f_s}{6}
			\quad\begin{cases}
				\text{audio:} & 44\,100/6 = 7350\ \text{Hz},\\[2pt]
				\text{vibración:} & 1102{,}5/6 = 183{,}75\ \text{Hz};
			\end{cases}
		\]
		\[
			\lambda=\pi
			\ \Longrightarrow\
			f = \frac{f_s}{2}
			\quad\begin{cases}
				\text{audio:} & 22{,}05\ \text{kHz},\\[2pt]
				\text{vibración:} & 551{,}25\ \text{Hz}.
			\end{cases}
		\]
		En los dos renglones la razón entre las frecuencias físicas es $40$, la misma que hay
		entre las frecuencias de muestreo. La explicación está en la definición: $\lambda$ es
		un \emph{ángulo por muestra}, y por sí solo no dice nada sobre segundos. Un avance de
		$\pi/3$ por muestra describe la misma secuencia en los dos casos; lo que cambia es
		cuánto tiempo pasa entre muestra y muestra, y ese dato lo aporta $\Delta t$. Para
		volver al mundo físico hay que reponer el $\Delta t$ que la definición de $\lambda$
		había absorbido.

	\end{enumerate}

	\begin{figure}[H]
		\centering
		\includestandalone[width=0.95\linewidth]{figures/fig_sol_ej18/fig_sol_ej18}
		\caption{Ejercicio~18. Las dos señales sobre el eje de frecuencia digital.
		(a) Audio muestreado a $44{,}1$ kHz. (b) Vibración muestreada a $1102{,}5$ Hz.
		El dibujo es el mismo en los dos paneles ---la banda ocupada, la banda base
		sombreada y las imágenes separadas $2\pi$---, y lo único que cambia es la regla
		en hertz debajo del eje. El extremo inferior de cada banda, $0{,}0009\pi$, no se
		distingue del origen a esta escala; el tramo negativo es la imagen de las
		frecuencias negativas.}
		\label{fig:sol_ej18}
	\end{figure}

\end{document}
